\chapter{Data Dictionary}\label{app:dict}
This appendix lists every table of the database with its columns, types, nullability, key
role and notes. It was produced from the live schema of the running database, not typed by
hand. ``bigint u'' means an unsigned 64-bit integer; PK, UK, FK and idx mean primary key,
unique key, foreign key and secondary index; \emph{generated} columns are computed by the
database.

\begingroup\footnotesize
\subsection*{Reference data (02\_reference\_tables)}
\begin{longtable}{@{}p{3.4cm}p{2.9cm}c c p{5.6cm}@{}}
\caption*{\textbf{regions}}\\
\toprule \textbf{Column} & \textbf{Type} & \textbf{Null} & \textbf{Key} & \textbf{Notes} \\ \midrule \endfirsthead
\toprule \textbf{Column} & \textbf{Type} & \textbf{Null} & \textbf{Key} & \textbf{Notes} \\ \midrule \endhead
\bottomrule \endlastfoot
region\_\allowbreak{}id & bigint u & N & PK & auto-increment \\
region\_\allowbreak{}name & varchar(100) & N &  &  \\
region\_\allowbreak{}code & varchar(10) & N & UK &  \\
parent\_\allowbreak{}region\_\allowbreak{}id & bigint u & Y & idx & FK to regions.region\_\allowbreak{}id, ON DELETE RESTRICT \\
\end{longtable}
\vspace{-0.4em}
\begin{longtable}{@{}p{3.4cm}p{2.9cm}c c p{5.6cm}@{}}
\caption*{\textbf{vehicle\_\allowbreak{}manufacturers}}\\
\toprule \textbf{Column} & \textbf{Type} & \textbf{Null} & \textbf{Key} & \textbf{Notes} \\ \midrule \endfirsthead
\toprule \textbf{Column} & \textbf{Type} & \textbf{Null} & \textbf{Key} & \textbf{Notes} \\ \midrule \endhead
\bottomrule \endlastfoot
manufacturer\_\allowbreak{}id & bigint u & N & PK & auto-increment \\
name & varchar(100) & N & UK &  \\
\end{longtable}
\vspace{-0.4em}
\begin{longtable}{@{}p{3.4cm}p{2.9cm}c c p{5.6cm}@{}}
\caption*{\textbf{vehicle\_\allowbreak{}models}}\\
\toprule \textbf{Column} & \textbf{Type} & \textbf{Null} & \textbf{Key} & \textbf{Notes} \\ \midrule \endfirsthead
\toprule \textbf{Column} & \textbf{Type} & \textbf{Null} & \textbf{Key} & \textbf{Notes} \\ \midrule \endhead
\bottomrule \endlastfoot
model\_\allowbreak{}id & bigint u & N & PK & auto-increment \\
manufacturer\_\allowbreak{}id & bigint u & N & idx & FK to vehicle\_\allowbreak{}manufacturers.manufacturer\_\allowbreak{}id, ON DELETE RESTRICT \\
model\_\allowbreak{}name & varchar(100) & N &  &  \\
\end{longtable}
\vspace{-0.4em}
\begin{longtable}{@{}p{3.4cm}p{2.9cm}c c p{5.6cm}@{}}
\caption*{\textbf{vehicle\_\allowbreak{}types}}\\
\toprule \textbf{Column} & \textbf{Type} & \textbf{Null} & \textbf{Key} & \textbf{Notes} \\ \midrule \endfirsthead
\toprule \textbf{Column} & \textbf{Type} & \textbf{Null} & \textbf{Key} & \textbf{Notes} \\ \midrule \endhead
\bottomrule \endlastfoot
vehicle\_\allowbreak{}type\_\allowbreak{}id & bigint u & N & PK & auto-increment \\
type\_\allowbreak{}name & varchar(50) & N & UK &  \\
is\_\allowbreak{}commercial & tinyint(1) & N &  & default 0 \\
\end{longtable}
\vspace{-0.4em}
\begin{longtable}{@{}p{3.4cm}p{2.9cm}c c p{5.6cm}@{}}
\caption*{\textbf{fuel\_\allowbreak{}types}}\\
\toprule \textbf{Column} & \textbf{Type} & \textbf{Null} & \textbf{Key} & \textbf{Notes} \\ \midrule \endfirsthead
\toprule \textbf{Column} & \textbf{Type} & \textbf{Null} & \textbf{Key} & \textbf{Notes} \\ \midrule \endhead
\bottomrule \endlastfoot
fuel\_\allowbreak{}type\_\allowbreak{}id & bigint u & N & PK & auto-increment \\
fuel\_\allowbreak{}name & varchar(30) & N & UK &  \\
\end{longtable}
\vspace{-0.4em}
\begin{longtable}{@{}p{3.4cm}p{2.9cm}c c p{5.6cm}@{}}
\caption*{\textbf{licence\_\allowbreak{}classes}}\\
\toprule \textbf{Column} & \textbf{Type} & \textbf{Null} & \textbf{Key} & \textbf{Notes} \\ \midrule \endfirsthead
\toprule \textbf{Column} & \textbf{Type} & \textbf{Null} & \textbf{Key} & \textbf{Notes} \\ \midrule \endhead
\bottomrule \endlastfoot
licence\_\allowbreak{}class\_\allowbreak{}id & bigint u & N & PK & auto-increment \\
class\_\allowbreak{}code & varchar(10) & N & UK &  \\
description & varchar(150) & N &  &  \\
\end{longtable}
\vspace{-0.4em}
\begin{longtable}{@{}p{3.4cm}p{2.9cm}c c p{5.6cm}@{}}
\caption*{\textbf{document\_\allowbreak{}types}}\\
\toprule \textbf{Column} & \textbf{Type} & \textbf{Null} & \textbf{Key} & \textbf{Notes} \\ \midrule \endfirsthead
\toprule \textbf{Column} & \textbf{Type} & \textbf{Null} & \textbf{Key} & \textbf{Notes} \\ \midrule \endhead
\bottomrule \endlastfoot
document\_\allowbreak{}type\_\allowbreak{}id & bigint u & N & PK & auto-increment \\
type\_\allowbreak{}name & varchar(100) & N & UK &  \\
is\_\allowbreak{}mandatory\_\allowbreak{}default & tinyint(1) & N &  & default 1 \\
validity\_\allowbreak{}period\_\allowbreak{}days & int u & Y &  &  \\
\end{longtable}
\vspace{-0.4em}
\begin{longtable}{@{}p{3.4cm}p{2.9cm}c c p{5.6cm}@{}}
\caption*{\textbf{service\_\allowbreak{}types}}\\
\toprule \textbf{Column} & \textbf{Type} & \textbf{Null} & \textbf{Key} & \textbf{Notes} \\ \midrule \endfirsthead
\toprule \textbf{Column} & \textbf{Type} & \textbf{Null} & \textbf{Key} & \textbf{Notes} \\ \midrule \endhead
\bottomrule \endlastfoot
service\_\allowbreak{}type\_\allowbreak{}id & bigint u & N & PK & auto-increment \\
service\_\allowbreak{}name & varchar(100) & N &  &  \\
service\_\allowbreak{}code & varchar(30) & N & UK &  \\
base\_\allowbreak{}fee & decimal(10,2) & N &  & default 0.00 \\
sla\_\allowbreak{}days & int u & N &  & default 7 \\
\end{longtable}
\vspace{-0.4em}
\begin{longtable}{@{}p{3.4cm}p{2.9cm}c c p{5.6cm}@{}}
\caption*{\textbf{violation\_\allowbreak{}types}}\\
\toprule \textbf{Column} & \textbf{Type} & \textbf{Null} & \textbf{Key} & \textbf{Notes} \\ \midrule \endfirsthead
\toprule \textbf{Column} & \textbf{Type} & \textbf{Null} & \textbf{Key} & \textbf{Notes} \\ \midrule \endhead
\bottomrule \endlastfoot
violation\_\allowbreak{}type\_\allowbreak{}id & bigint u & N & PK & auto-increment \\
description & varchar(150) & N & UK &  \\
base\_\allowbreak{}fine\_\allowbreak{}amount & decimal(10,2) & N &  &  \\
is\_\allowbreak{}cognizable & tinyint(1) & N &  & default 0 \\
\end{longtable}
\vspace{-0.4em}
\begin{longtable}{@{}p{3.4cm}p{2.9cm}c c p{5.6cm}@{}}
\caption*{\textbf{permit\_\allowbreak{}types}}\\
\toprule \textbf{Column} & \textbf{Type} & \textbf{Null} & \textbf{Key} & \textbf{Notes} \\ \midrule \endfirsthead
\toprule \textbf{Column} & \textbf{Type} & \textbf{Null} & \textbf{Key} & \textbf{Notes} \\ \midrule \endhead
\bottomrule \endlastfoot
permit\_\allowbreak{}type\_\allowbreak{}id & bigint u & N & PK & auto-increment \\
type\_\allowbreak{}name & varchar(50) & N & UK &  \\
validity\_\allowbreak{}months & int u & N &  &  \\
\end{longtable}
\vspace{-0.4em}
\begin{longtable}{@{}p{3.4cm}p{2.9cm}c c p{5.6cm}@{}}
\caption*{\textbf{payable\_\allowbreak{}types}}\\
\toprule \textbf{Column} & \textbf{Type} & \textbf{Null} & \textbf{Key} & \textbf{Notes} \\ \midrule \endfirsthead
\toprule \textbf{Column} & \textbf{Type} & \textbf{Null} & \textbf{Key} & \textbf{Notes} \\ \midrule \endhead
\bottomrule \endlastfoot
payable\_\allowbreak{}type\_\allowbreak{}id & bigint u & N & PK & auto-increment \\
type\_\allowbreak{}name & varchar(30) & N & UK &  \\
\end{longtable}
\vspace{-0.4em}
\subsection*{Identity (03\_identity)}
\begin{longtable}{@{}p{3.4cm}p{2.9cm}c c p{5.6cm}@{}}
\caption*{\textbf{persons}}\\
\toprule \textbf{Column} & \textbf{Type} & \textbf{Null} & \textbf{Key} & \textbf{Notes} \\ \midrule \endfirsthead
\toprule \textbf{Column} & \textbf{Type} & \textbf{Null} & \textbf{Key} & \textbf{Notes} \\ \midrule \endhead
\bottomrule \endlastfoot
person\_\allowbreak{}id & bigint u & N & PK & auto-increment \\
first\_\allowbreak{}name & varchar(60) & N &  &  \\
last\_\allowbreak{}name & varchar(60) & N & idx &  \\
date\_\allowbreak{}of\_\allowbreak{}birth & date & N &  &  \\
gender & enum('MALE',\allowbreak{}'FEMALE',\allowbreak{}'OTHER') & N &  &  \\
national\_\allowbreak{}id\_\allowbreak{}number & varchar(20) & N & UK &  \\
phone\_\allowbreak{}primary & varchar(15) & N & idx &  \\
phone\_\allowbreak{}secondary & varchar(15) & Y &  &  \\
email & varchar(120) & Y &  &  \\
photo\_\allowbreak{}path & varchar(255) & Y &  &  \\
created\_\allowbreak{}at & timestamp & N &  & default CURRENT\_\allowbreak{}TIMESTAMP \\
updated\_\allowbreak{}at & timestamp & N &  & default CURRENT\_\allowbreak{}TIMESTAMP \\
\end{longtable}
\vspace{-0.4em}
\begin{longtable}{@{}p{3.4cm}p{2.9cm}c c p{5.6cm}@{}}
\caption*{\textbf{addresses}}\\
\toprule \textbf{Column} & \textbf{Type} & \textbf{Null} & \textbf{Key} & \textbf{Notes} \\ \midrule \endfirsthead
\toprule \textbf{Column} & \textbf{Type} & \textbf{Null} & \textbf{Key} & \textbf{Notes} \\ \midrule \endhead
\bottomrule \endlastfoot
address\_\allowbreak{}id & bigint u & N & PK & auto-increment \\
person\_\allowbreak{}id & bigint u & N & idx & FK to persons.person\_\allowbreak{}id, ON DELETE CASCADE \\
line1 & varchar(150) & N &  &  \\
line2 & varchar(150) & Y &  &  \\
city & varchar(80) & N &  &  \\
state & varchar(80) & N &  &  \\
pincode & varchar(10) & N & idx &  \\
address\_\allowbreak{}type & enum('PERMANENT',\allowbreak{}'CURRENT',\allowbreak{}'OFFICE') & N &  & default CURRENT \\
is\_\allowbreak{}current & tinyint(1) & N &  & default 1 \\
valid\_\allowbreak{}from & date & N &  &  \\
valid\_\allowbreak{}to & date & Y &  &  \\
\end{longtable}
\vspace{-0.4em}
\begin{longtable}{@{}p{3.4cm}p{2.9cm}c c p{5.6cm}@{}}
\caption*{\textbf{citizens}}\\
\toprule \textbf{Column} & \textbf{Type} & \textbf{Null} & \textbf{Key} & \textbf{Notes} \\ \midrule \endfirsthead
\toprule \textbf{Column} & \textbf{Type} & \textbf{Null} & \textbf{Key} & \textbf{Notes} \\ \midrule \endhead
\bottomrule \endlastfoot
citizen\_\allowbreak{}id & bigint u & N & PK & auto-increment \\
person\_\allowbreak{}id & bigint u & N & UK & FK to persons.person\_\allowbreak{}id, ON DELETE RESTRICT \\
citizen\_\allowbreak{}code & varchar(20) & N & UK &  \\
blacklisted & tinyint(1) & N &  & default 0 \\
blacklist\_\allowbreak{}reason & varchar(255) & Y &  &  \\
registered\_\allowbreak{}at & timestamp & N &  & default CURRENT\_\allowbreak{}TIMESTAMP \\
\end{longtable}
\vspace{-0.4em}
\subsection*{Office structure (04\_office\_structure)}
\begin{longtable}{@{}p{3.4cm}p{2.9cm}c c p{5.6cm}@{}}
\caption*{\textbf{rto\_\allowbreak{}offices}}\\
\toprule \textbf{Column} & \textbf{Type} & \textbf{Null} & \textbf{Key} & \textbf{Notes} \\ \midrule \endfirsthead
\toprule \textbf{Column} & \textbf{Type} & \textbf{Null} & \textbf{Key} & \textbf{Notes} \\ \midrule \endhead
\bottomrule \endlastfoot
office\_\allowbreak{}id & bigint u & N & PK & auto-increment \\
region\_\allowbreak{}id & bigint u & N & idx & FK to regions.region\_\allowbreak{}id, ON DELETE RESTRICT \\
office\_\allowbreak{}name & varchar(100) & N &  &  \\
office\_\allowbreak{}code & varchar(15) & N & UK &  \\
address\_\allowbreak{}line & varchar(200) & N &  &  \\
is\_\allowbreak{}active & tinyint(1) & N &  & default 1 \\
\end{longtable}
\vspace{-0.4em}
\begin{longtable}{@{}p{3.4cm}p{2.9cm}c c p{5.6cm}@{}}
\caption*{\textbf{departments}}\\
\toprule \textbf{Column} & \textbf{Type} & \textbf{Null} & \textbf{Key} & \textbf{Notes} \\ \midrule \endfirsthead
\toprule \textbf{Column} & \textbf{Type} & \textbf{Null} & \textbf{Key} & \textbf{Notes} \\ \midrule \endhead
\bottomrule \endlastfoot
department\_\allowbreak{}id & bigint u & N & PK & auto-increment \\
department\_\allowbreak{}name & varchar(80) & N & UK &  \\
\end{longtable}
\vspace{-0.4em}
\begin{longtable}{@{}p{3.4cm}p{2.9cm}c c p{5.6cm}@{}}
\caption*{\textbf{designations}}\\
\toprule \textbf{Column} & \textbf{Type} & \textbf{Null} & \textbf{Key} & \textbf{Notes} \\ \midrule \endfirsthead
\toprule \textbf{Column} & \textbf{Type} & \textbf{Null} & \textbf{Key} & \textbf{Notes} \\ \midrule \endhead
\bottomrule \endlastfoot
designation\_\allowbreak{}id & bigint u & N & PK & auto-increment \\
title & varchar(80) & N & UK &  \\
rank\_\allowbreak{}level & tinyint u & N &  &  \\
\end{longtable}
\vspace{-0.4em}
\begin{longtable}{@{}p{3.4cm}p{2.9cm}c c p{5.6cm}@{}}
\caption*{\textbf{counters}}\\
\toprule \textbf{Column} & \textbf{Type} & \textbf{Null} & \textbf{Key} & \textbf{Notes} \\ \midrule \endfirsthead
\toprule \textbf{Column} & \textbf{Type} & \textbf{Null} & \textbf{Key} & \textbf{Notes} \\ \midrule \endhead
\bottomrule \endlastfoot
counter\_\allowbreak{}id & bigint u & N & PK & auto-increment \\
office\_\allowbreak{}id & bigint u & N & idx & FK to rto\_\allowbreak{}offices.office\_\allowbreak{}id, ON DELETE CASCADE \\
counter\_\allowbreak{}number & varchar(10) & N &  &  \\
service\_\allowbreak{}category & varchar(50) & Y &  &  \\
\end{longtable}
\vspace{-0.4em}
\subsection*{Employees (05\_employees)}
\begin{longtable}{@{}p{3.4cm}p{2.9cm}c c p{5.6cm}@{}}
\caption*{\textbf{employees}}\\
\toprule \textbf{Column} & \textbf{Type} & \textbf{Null} & \textbf{Key} & \textbf{Notes} \\ \midrule \endfirsthead
\toprule \textbf{Column} & \textbf{Type} & \textbf{Null} & \textbf{Key} & \textbf{Notes} \\ \midrule \endhead
\bottomrule \endlastfoot
employee\_\allowbreak{}id & bigint u & N & PK & auto-increment \\
person\_\allowbreak{}id & bigint u & N & UK & FK to persons.person\_\allowbreak{}id, ON DELETE RESTRICT \\
employee\_\allowbreak{}code & varchar(20) & N & UK &  \\
designation\_\allowbreak{}id & bigint u & N & idx & FK to designations.designation\_\allowbreak{}id, ON DELETE RESTRICT \\
date\_\allowbreak{}joined & date & N &  &  \\
is\_\allowbreak{}active & tinyint(1) & N &  & default 1 \\
\end{longtable}
\vspace{-0.4em}
\begin{longtable}{@{}p{3.4cm}p{2.9cm}c c p{5.6cm}@{}}
\caption*{\textbf{employee\_\allowbreak{}postings}}\\
\toprule \textbf{Column} & \textbf{Type} & \textbf{Null} & \textbf{Key} & \textbf{Notes} \\ \midrule \endfirsthead
\toprule \textbf{Column} & \textbf{Type} & \textbf{Null} & \textbf{Key} & \textbf{Notes} \\ \midrule \endhead
\bottomrule \endlastfoot
posting\_\allowbreak{}id & bigint u & N & PK & auto-increment \\
employee\_\allowbreak{}id & bigint u & N & idx & FK to employees.employee\_\allowbreak{}id, ON DELETE CASCADE \\
office\_\allowbreak{}id & bigint u & N & idx & FK to rto\_\allowbreak{}offices.office\_\allowbreak{}id, ON DELETE RESTRICT \\
department\_\allowbreak{}id & bigint u & N & idx & FK to departments.department\_\allowbreak{}id, ON DELETE RESTRICT \\
posted\_\allowbreak{}from & date & N &  &  \\
posted\_\allowbreak{}to & date & Y &  &  \\
current\_\allowbreak{}flag & tinyint & Y &  & generated: if((posted\_\allowbreak{}to is null),1,NULL) \\
\end{longtable}
\vspace{-0.4em}
\subsection*{Access control (06\_rbac)}
\begin{longtable}{@{}p{3.4cm}p{2.9cm}c c p{5.6cm}@{}}
\caption*{\textbf{users}}\\
\toprule \textbf{Column} & \textbf{Type} & \textbf{Null} & \textbf{Key} & \textbf{Notes} \\ \midrule \endfirsthead
\toprule \textbf{Column} & \textbf{Type} & \textbf{Null} & \textbf{Key} & \textbf{Notes} \\ \midrule \endhead
\bottomrule \endlastfoot
user\_\allowbreak{}id & bigint u & N & PK & auto-increment \\
person\_\allowbreak{}id & bigint u & N & idx & FK to persons.person\_\allowbreak{}id, ON DELETE CASCADE \\
username & varchar(50) & N & UK &  \\
password\_\allowbreak{}hash & varchar(255) & N &  &  \\
is\_\allowbreak{}active & tinyint(1) & N &  & default 1 \\
last\_\allowbreak{}login\_\allowbreak{}at & datetime & Y &  &  \\
\end{longtable}
\vspace{-0.4em}
\begin{longtable}{@{}p{3.4cm}p{2.9cm}c c p{5.6cm}@{}}
\caption*{\textbf{roles}}\\
\toprule \textbf{Column} & \textbf{Type} & \textbf{Null} & \textbf{Key} & \textbf{Notes} \\ \midrule \endfirsthead
\toprule \textbf{Column} & \textbf{Type} & \textbf{Null} & \textbf{Key} & \textbf{Notes} \\ \midrule \endhead
\bottomrule \endlastfoot
role\_\allowbreak{}id & bigint u & N & PK & auto-increment \\
role\_\allowbreak{}name & varchar(50) & N & UK &  \\
\end{longtable}
\vspace{-0.4em}
\begin{longtable}{@{}p{3.4cm}p{2.9cm}c c p{5.6cm}@{}}
\caption*{\textbf{permissions}}\\
\toprule \textbf{Column} & \textbf{Type} & \textbf{Null} & \textbf{Key} & \textbf{Notes} \\ \midrule \endfirsthead
\toprule \textbf{Column} & \textbf{Type} & \textbf{Null} & \textbf{Key} & \textbf{Notes} \\ \midrule \endhead
\bottomrule \endlastfoot
permission\_\allowbreak{}id & bigint u & N & PK & auto-increment \\
permission\_\allowbreak{}key & varchar(80) & N & UK &  \\
\end{longtable}
\vspace{-0.4em}
\begin{longtable}{@{}p{3.4cm}p{2.9cm}c c p{5.6cm}@{}}
\caption*{\textbf{user\_\allowbreak{}roles}}\\
\toprule \textbf{Column} & \textbf{Type} & \textbf{Null} & \textbf{Key} & \textbf{Notes} \\ \midrule \endfirsthead
\toprule \textbf{Column} & \textbf{Type} & \textbf{Null} & \textbf{Key} & \textbf{Notes} \\ \midrule \endhead
\bottomrule \endlastfoot
user\_\allowbreak{}id & bigint u & N & PK & FK to users.user\_\allowbreak{}id, ON DELETE CASCADE \\
role\_\allowbreak{}id & bigint u & N & PK & FK to roles.role\_\allowbreak{}id, ON DELETE CASCADE \\
\end{longtable}
\vspace{-0.4em}
\begin{longtable}{@{}p{3.4cm}p{2.9cm}c c p{5.6cm}@{}}
\caption*{\textbf{role\_\allowbreak{}permissions}}\\
\toprule \textbf{Column} & \textbf{Type} & \textbf{Null} & \textbf{Key} & \textbf{Notes} \\ \midrule \endfirsthead
\toprule \textbf{Column} & \textbf{Type} & \textbf{Null} & \textbf{Key} & \textbf{Notes} \\ \midrule \endhead
\bottomrule \endlastfoot
role\_\allowbreak{}id & bigint u & N & PK & FK to roles.role\_\allowbreak{}id, ON DELETE CASCADE \\
permission\_\allowbreak{}id & bigint u & N & PK & FK to permissions.permission\_\allowbreak{}id, ON DELETE CASCADE \\
\end{longtable}
\vspace{-0.4em}
\subsection*{Applications (07\_applications)}
\begin{longtable}{@{}p{3.4cm}p{2.9cm}c c p{5.6cm}@{}}
\caption*{\textbf{applicants}}\\
\toprule \textbf{Column} & \textbf{Type} & \textbf{Null} & \textbf{Key} & \textbf{Notes} \\ \midrule \endfirsthead
\toprule \textbf{Column} & \textbf{Type} & \textbf{Null} & \textbf{Key} & \textbf{Notes} \\ \midrule \endhead
\bottomrule \endlastfoot
applicant\_\allowbreak{}id & bigint u & N & PK & auto-increment \\
citizen\_\allowbreak{}id & bigint u & N & UK & FK to citizens.citizen\_\allowbreak{}id, ON DELETE RESTRICT \\
preferred\_\allowbreak{}office\_\allowbreak{}id & bigint u & Y & idx & FK to rto\_\allowbreak{}offices.office\_\allowbreak{}id, ON DELETE SET NULL \\
\end{longtable}
\vspace{-0.4em}
\begin{longtable}{@{}p{3.4cm}p{2.9cm}c c p{5.6cm}@{}}
\caption*{\textbf{applications}}\\
\toprule \textbf{Column} & \textbf{Type} & \textbf{Null} & \textbf{Key} & \textbf{Notes} \\ \midrule \endfirsthead
\toprule \textbf{Column} & \textbf{Type} & \textbf{Null} & \textbf{Key} & \textbf{Notes} \\ \midrule \endhead
\bottomrule \endlastfoot
application\_\allowbreak{}id & bigint u & N & PK & auto-increment \\
application\_\allowbreak{}number & varchar(30) & N & UK &  \\
applicant\_\allowbreak{}id & bigint u & N & idx & FK to applicants.applicant\_\allowbreak{}id, ON DELETE RESTRICT \\
service\_\allowbreak{}type\_\allowbreak{}id & bigint u & N & idx & FK to service\_\allowbreak{}types.service\_\allowbreak{}type\_\allowbreak{}id, ON DELETE RESTRICT \\
office\_\allowbreak{}id & bigint u & N & idx & FK to rto\_\allowbreak{}offices.office\_\allowbreak{}id, ON DELETE RESTRICT \\
assigned\_\allowbreak{}officer\_\allowbreak{}id & bigint u & Y & idx & FK to employees.employee\_\allowbreak{}id, ON DELETE SET NULL \\
current\_\allowbreak{}status & enum('SUBMITTED',\allowbreak{}'DOCS\_\allowbreak{}PENDING',\allowbreak{}'UNDER\_\allowbreak{}VERIFICATION',\allowbreak{}'APPOINTMENT\_\allowbreak{}SCHEDULED',\allowbreak{}'AWAITING\_\allowbreak{}PAYMENT',\allowbreak{}'APPROVED',\allowbreak{}'REJECTED',\allowbreak{}'COMPLETED',\allowbreak{}'CANCELLED') & N & idx & default SUBMITTED \\
submitted\_\allowbreak{}at & datetime & N &  & default CURRENT\_\allowbreak{}TIMESTAMP \\
completed\_\allowbreak{}at & datetime & Y &  &  \\
remarks & varchar(500) & Y &  &  \\
\end{longtable}
\vspace{-0.4em}
\begin{longtable}{@{}p{3.4cm}p{2.9cm}c c p{5.6cm}@{}}
\caption*{\textbf{application\_\allowbreak{}status\_\allowbreak{}history}}\\
\toprule \textbf{Column} & \textbf{Type} & \textbf{Null} & \textbf{Key} & \textbf{Notes} \\ \midrule \endfirsthead
\toprule \textbf{Column} & \textbf{Type} & \textbf{Null} & \textbf{Key} & \textbf{Notes} \\ \midrule \endhead
\bottomrule \endlastfoot
history\_\allowbreak{}id & bigint u & N & PK & auto-increment \\
application\_\allowbreak{}id & bigint u & N & idx & FK to applications.application\_\allowbreak{}id, ON DELETE CASCADE \\
previous\_\allowbreak{}status & varchar(30) & Y &  &  \\
new\_\allowbreak{}status & varchar(30) & N &  &  \\
changed\_\allowbreak{}by\_\allowbreak{}user\_\allowbreak{}id & bigint u & Y & idx & FK to users.user\_\allowbreak{}id, ON DELETE SET NULL \\
changed\_\allowbreak{}at & timestamp & N &  & default CURRENT\_\allowbreak{}TIMESTAMP \\
reason & varchar(255) & Y &  &  \\
\end{longtable}
\vspace{-0.4em}
\begin{longtable}{@{}p{3.4cm}p{2.9cm}c c p{5.6cm}@{}}
\caption*{\textbf{documents}}\\
\toprule \textbf{Column} & \textbf{Type} & \textbf{Null} & \textbf{Key} & \textbf{Notes} \\ \midrule \endfirsthead
\toprule \textbf{Column} & \textbf{Type} & \textbf{Null} & \textbf{Key} & \textbf{Notes} \\ \midrule \endhead
\bottomrule \endlastfoot
document\_\allowbreak{}id & bigint u & N & PK & auto-increment \\
application\_\allowbreak{}id & bigint u & N & idx & FK to applications.application\_\allowbreak{}id, ON DELETE CASCADE \\
document\_\allowbreak{}type\_\allowbreak{}id & bigint u & N & idx & FK to document\_\allowbreak{}types.document\_\allowbreak{}type\_\allowbreak{}id, ON DELETE RESTRICT \\
file\_\allowbreak{}path & varchar(255) & N &  &  \\
verification\_\allowbreak{}status & enum('PENDING',\allowbreak{}'VERIFIED',\allowbreak{}'REJECTED',\allowbreak{}'EXPIRED') & N &  & default PENDING \\
verified\_\allowbreak{}by\_\allowbreak{}employee\_\allowbreak{}id & bigint u & Y & idx & FK to employees.employee\_\allowbreak{}id, ON DELETE SET NULL \\
rejection\_\allowbreak{}reason & varchar(255) & Y &  &  \\
uploaded\_\allowbreak{}at & datetime & N &  & default CURRENT\_\allowbreak{}TIMESTAMP \\
verified\_\allowbreak{}at & datetime & Y &  &  \\
\end{longtable}
\vspace{-0.4em}
\begin{longtable}{@{}p{3.4cm}p{2.9cm}c c p{5.6cm}@{}}
\caption*{\textbf{appointment\_\allowbreak{}slots}}\\
\toprule \textbf{Column} & \textbf{Type} & \textbf{Null} & \textbf{Key} & \textbf{Notes} \\ \midrule \endfirsthead
\toprule \textbf{Column} & \textbf{Type} & \textbf{Null} & \textbf{Key} & \textbf{Notes} \\ \midrule \endhead
\bottomrule \endlastfoot
slot\_\allowbreak{}id & bigint u & N & PK & auto-increment \\
office\_\allowbreak{}id & bigint u & N & idx & FK to rto\_\allowbreak{}offices.office\_\allowbreak{}id, ON DELETE CASCADE \\
counter\_\allowbreak{}id & bigint u & Y & idx & FK to counters.counter\_\allowbreak{}id, ON DELETE SET NULL \\
slot\_\allowbreak{}date & date & N &  &  \\
start\_\allowbreak{}time & time & N &  &  \\
end\_\allowbreak{}time & time & N &  &  \\
capacity & smallint u & N &  & default 1 \\
booked\_\allowbreak{}count & smallint u & N &  & default 0 \\
\end{longtable}
\vspace{-0.4em}
\begin{longtable}{@{}p{3.4cm}p{2.9cm}c c p{5.6cm}@{}}
\caption*{\textbf{appointments}}\\
\toprule \textbf{Column} & \textbf{Type} & \textbf{Null} & \textbf{Key} & \textbf{Notes} \\ \midrule \endfirsthead
\toprule \textbf{Column} & \textbf{Type} & \textbf{Null} & \textbf{Key} & \textbf{Notes} \\ \midrule \endhead
\bottomrule \endlastfoot
appointment\_\allowbreak{}id & bigint u & N & PK & auto-increment \\
application\_\allowbreak{}id & bigint u & N & UK & FK to applications.application\_\allowbreak{}id, ON DELETE CASCADE \\
slot\_\allowbreak{}id & bigint u & N & idx & FK to appointment\_\allowbreak{}slots.slot\_\allowbreak{}id, ON DELETE RESTRICT \\
token\_\allowbreak{}number & varchar(15) & N &  &  \\
status & enum('BOOKED',\allowbreak{}'RESCHEDULED',\allowbreak{}'CANCELLED',\allowbreak{}'NO\_\allowbreak{}SHOW',\allowbreak{}'COMPLETED') & N &  & default BOOKED \\
booked\_\allowbreak{}at & datetime & N &  & default CURRENT\_\allowbreak{}TIMESTAMP \\
\end{longtable}
\vspace{-0.4em}
\subsection*{Licensing (08\_licence)}
\begin{longtable}{@{}p{3.4cm}p{2.9cm}c c p{5.6cm}@{}}
\caption*{\textbf{driving\_\allowbreak{}schools}}\\
\toprule \textbf{Column} & \textbf{Type} & \textbf{Null} & \textbf{Key} & \textbf{Notes} \\ \midrule \endfirsthead
\toprule \textbf{Column} & \textbf{Type} & \textbf{Null} & \textbf{Key} & \textbf{Notes} \\ \midrule \endhead
\bottomrule \endlastfoot
school\_\allowbreak{}id & bigint u & N & PK & auto-increment \\
school\_\allowbreak{}name & varchar(120) & N &  &  \\
licence\_\allowbreak{}number & varchar(30) & N & UK &  \\
office\_\allowbreak{}id & bigint u & N & idx & FK to rto\_\allowbreak{}offices.office\_\allowbreak{}id, ON DELETE RESTRICT \\
\end{longtable}
\vspace{-0.4em}
\begin{longtable}{@{}p{3.4cm}p{2.9cm}c c p{5.6cm}@{}}
\caption*{\textbf{driving\_\allowbreak{}instructors}}\\
\toprule \textbf{Column} & \textbf{Type} & \textbf{Null} & \textbf{Key} & \textbf{Notes} \\ \midrule \endfirsthead
\toprule \textbf{Column} & \textbf{Type} & \textbf{Null} & \textbf{Key} & \textbf{Notes} \\ \midrule \endhead
\bottomrule \endlastfoot
instructor\_\allowbreak{}id & bigint u & N & PK & auto-increment \\
person\_\allowbreak{}id & bigint u & N & UK & FK to persons.person\_\allowbreak{}id, ON DELETE RESTRICT \\
school\_\allowbreak{}id & bigint u & N & idx & FK to driving\_\allowbreak{}schools.school\_\allowbreak{}id, ON DELETE CASCADE \\
licence\_\allowbreak{}class\_\allowbreak{}id & bigint u & N & idx & FK to licence\_\allowbreak{}classes.licence\_\allowbreak{}class\_\allowbreak{}id, ON DELETE RESTRICT \\
\end{longtable}
\vspace{-0.4em}
\begin{longtable}{@{}p{3.4cm}p{2.9cm}c c p{5.6cm}@{}}
\caption*{\textbf{test\_\allowbreak{}centres}}\\
\toprule \textbf{Column} & \textbf{Type} & \textbf{Null} & \textbf{Key} & \textbf{Notes} \\ \midrule \endfirsthead
\toprule \textbf{Column} & \textbf{Type} & \textbf{Null} & \textbf{Key} & \textbf{Notes} \\ \midrule \endhead
\bottomrule \endlastfoot
test\_\allowbreak{}centre\_\allowbreak{}id & bigint u & N & PK & auto-increment \\
office\_\allowbreak{}id & bigint u & N & idx & FK to rto\_\allowbreak{}offices.office\_\allowbreak{}id, ON DELETE CASCADE \\
centre\_\allowbreak{}name & varchar(120) & N &  &  \\
\end{longtable}
\vspace{-0.4em}
\begin{longtable}{@{}p{3.4cm}p{2.9cm}c c p{5.6cm}@{}}
\caption*{\textbf{learner\_\allowbreak{}licences}}\\
\toprule \textbf{Column} & \textbf{Type} & \textbf{Null} & \textbf{Key} & \textbf{Notes} \\ \midrule \endfirsthead
\toprule \textbf{Column} & \textbf{Type} & \textbf{Null} & \textbf{Key} & \textbf{Notes} \\ \midrule \endhead
\bottomrule \endlastfoot
learner\_\allowbreak{}licence\_\allowbreak{}id & bigint u & N & PK & auto-increment \\
licence\_\allowbreak{}number & varchar(30) & N & UK &  \\
citizen\_\allowbreak{}id & bigint u & N & idx & FK to citizens.citizen\_\allowbreak{}id, ON DELETE RESTRICT \\
application\_\allowbreak{}id & bigint u & N & idx & FK to applications.application\_\allowbreak{}id, ON DELETE RESTRICT \\
office\_\allowbreak{}id & bigint u & N & idx & FK to rto\_\allowbreak{}offices.office\_\allowbreak{}id, ON DELETE RESTRICT \\
issue\_\allowbreak{}date & date & N &  &  \\
expiry\_\allowbreak{}date & date & N &  &  \\
status & enum('ACTIVE',\allowbreak{}'EXPIRED',\allowbreak{}'CONVERTED',\allowbreak{}'CANCELLED') & N &  & default ACTIVE \\
active\_\allowbreak{}flag & tinyint & Y &  & generated: if((status = \_\allowbreak{}utf8mb4'ACTIVE'),1,NULL) \\
\end{longtable}
\vspace{-0.4em}
\begin{longtable}{@{}p{3.4cm}p{2.9cm}c c p{5.6cm}@{}}
\caption*{\textbf{driving\_\allowbreak{}licences}}\\
\toprule \textbf{Column} & \textbf{Type} & \textbf{Null} & \textbf{Key} & \textbf{Notes} \\ \midrule \endfirsthead
\toprule \textbf{Column} & \textbf{Type} & \textbf{Null} & \textbf{Key} & \textbf{Notes} \\ \midrule \endhead
\bottomrule \endlastfoot
driving\_\allowbreak{}licence\_\allowbreak{}id & bigint u & N & PK & auto-increment \\
licence\_\allowbreak{}number & varchar(30) & N & UK &  \\
citizen\_\allowbreak{}id & bigint u & N & idx & FK to citizens.citizen\_\allowbreak{}id, ON DELETE RESTRICT \\
learner\_\allowbreak{}licence\_\allowbreak{}id & bigint u & Y & idx & FK to learner\_\allowbreak{}licences.learner\_\allowbreak{}licence\_\allowbreak{}id, ON DELETE SET NULL \\
application\_\allowbreak{}id & bigint u & N & idx & FK to applications.application\_\allowbreak{}id, ON DELETE RESTRICT \\
office\_\allowbreak{}id & bigint u & N & idx & FK to rto\_\allowbreak{}offices.office\_\allowbreak{}id, ON DELETE RESTRICT \\
issue\_\allowbreak{}date & date & N &  &  \\
expiry\_\allowbreak{}date & date & N &  &  \\
current\_\allowbreak{}status & enum('ACTIVE',\allowbreak{}'SUSPENDED',\allowbreak{}'REVOKED',\allowbreak{}'EXPIRED') & N & idx & default ACTIVE \\
\end{longtable}
\vspace{-0.4em}
\begin{longtable}{@{}p{3.4cm}p{2.9cm}c c p{5.6cm}@{}}
\caption*{\textbf{licence\_\allowbreak{}class\_\allowbreak{}assignments}}\\
\toprule \textbf{Column} & \textbf{Type} & \textbf{Null} & \textbf{Key} & \textbf{Notes} \\ \midrule \endfirsthead
\toprule \textbf{Column} & \textbf{Type} & \textbf{Null} & \textbf{Key} & \textbf{Notes} \\ \midrule \endhead
\bottomrule \endlastfoot
driving\_\allowbreak{}licence\_\allowbreak{}id & bigint u & N & PK & FK to driving\_\allowbreak{}licences.driving\_\allowbreak{}licence\_\allowbreak{}id, ON DELETE CASCADE \\
licence\_\allowbreak{}class\_\allowbreak{}id & bigint u & N & PK & FK to licence\_\allowbreak{}classes.licence\_\allowbreak{}class\_\allowbreak{}id, ON DELETE RESTRICT \\
granted\_\allowbreak{}on & date & N &  &  \\
\end{longtable}
\vspace{-0.4em}
\begin{longtable}{@{}p{3.4cm}p{2.9cm}c c p{5.6cm}@{}}
\caption*{\textbf{licence\_\allowbreak{}status\_\allowbreak{}history}}\\
\toprule \textbf{Column} & \textbf{Type} & \textbf{Null} & \textbf{Key} & \textbf{Notes} \\ \midrule \endfirsthead
\toprule \textbf{Column} & \textbf{Type} & \textbf{Null} & \textbf{Key} & \textbf{Notes} \\ \midrule \endhead
\bottomrule \endlastfoot
history\_\allowbreak{}id & bigint u & N & PK & auto-increment \\
driving\_\allowbreak{}licence\_\allowbreak{}id & bigint u & N & idx & FK to driving\_\allowbreak{}licences.driving\_\allowbreak{}licence\_\allowbreak{}id, ON DELETE CASCADE \\
previous\_\allowbreak{}status & varchar(20) & Y &  &  \\
new\_\allowbreak{}status & varchar(20) & N &  &  \\
changed\_\allowbreak{}by\_\allowbreak{}user\_\allowbreak{}id & bigint u & Y & idx & FK to users.user\_\allowbreak{}id, ON DELETE SET NULL \\
changed\_\allowbreak{}at & timestamp & N &  & default CURRENT\_\allowbreak{}TIMESTAMP \\
reason & varchar(255) & Y &  &  \\
\end{longtable}
\vspace{-0.4em}
\begin{longtable}{@{}p{3.4cm}p{2.9cm}c c p{5.6cm}@{}}
\caption*{\textbf{driving\_\allowbreak{}tests}}\\
\toprule \textbf{Column} & \textbf{Type} & \textbf{Null} & \textbf{Key} & \textbf{Notes} \\ \midrule \endfirsthead
\toprule \textbf{Column} & \textbf{Type} & \textbf{Null} & \textbf{Key} & \textbf{Notes} \\ \midrule \endhead
\bottomrule \endlastfoot
test\_\allowbreak{}id & bigint u & N & PK & auto-increment \\
application\_\allowbreak{}id & bigint u & N & idx & FK to applications.application\_\allowbreak{}id, ON DELETE CASCADE \\
citizen\_\allowbreak{}id & bigint u & N & idx & FK to citizens.citizen\_\allowbreak{}id, ON DELETE RESTRICT \\
test\_\allowbreak{}centre\_\allowbreak{}id & bigint u & N & idx & FK to test\_\allowbreak{}centres.test\_\allowbreak{}centre\_\allowbreak{}id, ON DELETE RESTRICT \\
examiner\_\allowbreak{}employee\_\allowbreak{}id & bigint u & N & idx & FK to employees.employee\_\allowbreak{}id, ON DELETE RESTRICT \\
scheduled\_\allowbreak{}at & datetime & N &  &  \\
result & enum('PENDING',\allowbreak{}'PASS',\allowbreak{}'FAIL',\allowbreak{}'ABSENT') & N &  & default PENDING \\
remarks & varchar(255) & Y &  &  \\
\end{longtable}
\vspace{-0.4em}
\subsection*{Vehicles (09\_vehicles)}
\begin{longtable}{@{}p{3.4cm}p{2.9cm}c c p{5.6cm}@{}}
\caption*{\textbf{vehicles}}\\
\toprule \textbf{Column} & \textbf{Type} & \textbf{Null} & \textbf{Key} & \textbf{Notes} \\ \midrule \endfirsthead
\toprule \textbf{Column} & \textbf{Type} & \textbf{Null} & \textbf{Key} & \textbf{Notes} \\ \midrule \endhead
\bottomrule \endlastfoot
vehicle\_\allowbreak{}id & bigint u & N & PK & auto-increment \\
registration\_\allowbreak{}number & varchar(20) & N & UK &  \\
chassis\_\allowbreak{}number & varchar(40) & N & UK &  \\
engine\_\allowbreak{}number & varchar(40) & N & UK &  \\
manufacturer\_\allowbreak{}id & bigint u & N & idx & FK to vehicle\_\allowbreak{}manufacturers.manufacturer\_\allowbreak{}id, ON DELETE RESTRICT \\
model\_\allowbreak{}id & bigint u & N & idx & FK to vehicle\_\allowbreak{}models.model\_\allowbreak{}id, ON DELETE RESTRICT \\
vehicle\_\allowbreak{}type\_\allowbreak{}id & bigint u & N & idx & FK to vehicle\_\allowbreak{}types.vehicle\_\allowbreak{}type\_\allowbreak{}id, ON DELETE RESTRICT \\
fuel\_\allowbreak{}type\_\allowbreak{}id & bigint u & N & idx & FK to fuel\_\allowbreak{}types.fuel\_\allowbreak{}type\_\allowbreak{}id, ON DELETE RESTRICT \\
manufacture\_\allowbreak{}year & year & N &  &  \\
color & varchar(30) & Y &  &  \\
registering\_\allowbreak{}office\_\allowbreak{}id & bigint u & N & idx & FK to rto\_\allowbreak{}offices.office\_\allowbreak{}id, ON DELETE RESTRICT \\
registration\_\allowbreak{}date & date & N &  &  \\
status & enum('ACTIVE',\allowbreak{}'BLACKLISTED',\allowbreak{}'SCRAPPED',\allowbreak{}'DEREGISTERED') & N &  & default ACTIVE \\
\end{longtable}
\vspace{-0.4em}
\begin{longtable}{@{}p{3.4cm}p{2.9cm}c c p{5.6cm}@{}}
\caption*{\textbf{vehicle\_\allowbreak{}ownerships}}\\
\toprule \textbf{Column} & \textbf{Type} & \textbf{Null} & \textbf{Key} & \textbf{Notes} \\ \midrule \endfirsthead
\toprule \textbf{Column} & \textbf{Type} & \textbf{Null} & \textbf{Key} & \textbf{Notes} \\ \midrule \endhead
\bottomrule \endlastfoot
ownership\_\allowbreak{}id & bigint u & N & PK & auto-increment \\
vehicle\_\allowbreak{}id & bigint u & N & idx & FK to vehicles.vehicle\_\allowbreak{}id, ON DELETE RESTRICT \\
citizen\_\allowbreak{}id & bigint u & N & idx & FK to citizens.citizen\_\allowbreak{}id, ON DELETE RESTRICT \\
effective\_\allowbreak{}from & date & N &  &  \\
effective\_\allowbreak{}to & date & Y &  &  \\
current\_\allowbreak{}flag & tinyint & Y &  & generated: if((effective\_\allowbreak{}to is null),1,NULL) \\
\end{longtable}
\vspace{-0.4em}
\begin{longtable}{@{}p{3.4cm}p{2.9cm}c c p{5.6cm}@{}}
\caption*{\textbf{ownership\_\allowbreak{}transfers}}\\
\toprule \textbf{Column} & \textbf{Type} & \textbf{Null} & \textbf{Key} & \textbf{Notes} \\ \midrule \endfirsthead
\toprule \textbf{Column} & \textbf{Type} & \textbf{Null} & \textbf{Key} & \textbf{Notes} \\ \midrule \endhead
\bottomrule \endlastfoot
transfer\_\allowbreak{}id & bigint u & N & PK & auto-increment \\
vehicle\_\allowbreak{}id & bigint u & N & idx & FK to vehicles.vehicle\_\allowbreak{}id, ON DELETE RESTRICT \\
from\_\allowbreak{}citizen\_\allowbreak{}id & bigint u & N & idx & FK to citizens.citizen\_\allowbreak{}id, ON DELETE RESTRICT \\
to\_\allowbreak{}citizen\_\allowbreak{}id & bigint u & N & idx & FK to citizens.citizen\_\allowbreak{}id, ON DELETE RESTRICT \\
application\_\allowbreak{}id & bigint u & N & idx & FK to applications.application\_\allowbreak{}id, ON DELETE RESTRICT \\
requested\_\allowbreak{}at & datetime & N &  & default CURRENT\_\allowbreak{}TIMESTAMP \\
approved\_\allowbreak{}at & datetime & Y &  &  \\
status & enum('PENDING',\allowbreak{}'APPROVED',\allowbreak{}'REJECTED') & N &  & default PENDING \\
\end{longtable}
\vspace{-0.4em}
\subsection*{Compliance (10\_compliance)}
\begin{longtable}{@{}p{3.4cm}p{2.9cm}c c p{5.6cm}@{}}
\caption*{\textbf{vehicle\_\allowbreak{}inspections}}\\
\toprule \textbf{Column} & \textbf{Type} & \textbf{Null} & \textbf{Key} & \textbf{Notes} \\ \midrule \endfirsthead
\toprule \textbf{Column} & \textbf{Type} & \textbf{Null} & \textbf{Key} & \textbf{Notes} \\ \midrule \endhead
\bottomrule \endlastfoot
inspection\_\allowbreak{}id & bigint u & N & PK & auto-increment \\
vehicle\_\allowbreak{}id & bigint u & N & idx & FK to vehicles.vehicle\_\allowbreak{}id, ON DELETE CASCADE \\
inspector\_\allowbreak{}employee\_\allowbreak{}id & bigint u & N & idx & FK to employees.employee\_\allowbreak{}id, ON DELETE RESTRICT \\
inspected\_\allowbreak{}at & datetime & N &  &  \\
result & enum('PASS',\allowbreak{}'FAIL') & N &  &  \\
remarks & varchar(255) & Y &  &  \\
\end{longtable}
\vspace{-0.4em}
\begin{longtable}{@{}p{3.4cm}p{2.9cm}c c p{5.6cm}@{}}
\caption*{\textbf{fitness\_\allowbreak{}certificates}}\\
\toprule \textbf{Column} & \textbf{Type} & \textbf{Null} & \textbf{Key} & \textbf{Notes} \\ \midrule \endfirsthead
\toprule \textbf{Column} & \textbf{Type} & \textbf{Null} & \textbf{Key} & \textbf{Notes} \\ \midrule \endhead
\bottomrule \endlastfoot
certificate\_\allowbreak{}id & bigint u & N & PK & auto-increment \\
vehicle\_\allowbreak{}id & bigint u & N & idx & FK to vehicles.vehicle\_\allowbreak{}id, ON DELETE CASCADE \\
inspection\_\allowbreak{}id & bigint u & Y & idx & FK to vehicle\_\allowbreak{}inspections.inspection\_\allowbreak{}id, ON DELETE SET NULL \\
issue\_\allowbreak{}date & date & N &  &  \\
expiry\_\allowbreak{}date & date & N &  &  \\
status & enum('ACTIVE',\allowbreak{}'EXPIRED',\allowbreak{}'REVOKED') & N &  & default ACTIVE \\
\end{longtable}
\vspace{-0.4em}
\begin{longtable}{@{}p{3.4cm}p{2.9cm}c c p{5.6cm}@{}}
\caption*{\textbf{pollution\_\allowbreak{}certificates}}\\
\toprule \textbf{Column} & \textbf{Type} & \textbf{Null} & \textbf{Key} & \textbf{Notes} \\ \midrule \endfirsthead
\toprule \textbf{Column} & \textbf{Type} & \textbf{Null} & \textbf{Key} & \textbf{Notes} \\ \midrule \endhead
\bottomrule \endlastfoot
puc\_\allowbreak{}id & bigint u & N & PK & auto-increment \\
vehicle\_\allowbreak{}id & bigint u & N & idx & FK to vehicles.vehicle\_\allowbreak{}id, ON DELETE CASCADE \\
issue\_\allowbreak{}date & date & N &  &  \\
expiry\_\allowbreak{}date & date & N &  &  \\
status & enum('ACTIVE',\allowbreak{}'EXPIRED') & N &  & default ACTIVE \\
\end{longtable}
\vspace{-0.4em}
\begin{longtable}{@{}p{3.4cm}p{2.9cm}c c p{5.6cm}@{}}
\caption*{\textbf{insurance\_\allowbreak{}policies}}\\
\toprule \textbf{Column} & \textbf{Type} & \textbf{Null} & \textbf{Key} & \textbf{Notes} \\ \midrule \endfirsthead
\toprule \textbf{Column} & \textbf{Type} & \textbf{Null} & \textbf{Key} & \textbf{Notes} \\ \midrule \endhead
\bottomrule \endlastfoot
policy\_\allowbreak{}id & bigint u & N & PK & auto-increment \\
vehicle\_\allowbreak{}id & bigint u & N & idx & FK to vehicles.vehicle\_\allowbreak{}id, ON DELETE CASCADE \\
provider\_\allowbreak{}name & varchar(100) & N &  &  \\
policy\_\allowbreak{}number & varchar(40) & N & UK &  \\
start\_\allowbreak{}date & date & N &  &  \\
end\_\allowbreak{}date & date & N &  &  \\
\end{longtable}
\vspace{-0.4em}
\begin{longtable}{@{}p{3.4cm}p{2.9cm}c c p{5.6cm}@{}}
\caption*{\textbf{road\_\allowbreak{}tax\_\allowbreak{}records}}\\
\toprule \textbf{Column} & \textbf{Type} & \textbf{Null} & \textbf{Key} & \textbf{Notes} \\ \midrule \endfirsthead
\toprule \textbf{Column} & \textbf{Type} & \textbf{Null} & \textbf{Key} & \textbf{Notes} \\ \midrule \endhead
\bottomrule \endlastfoot
tax\_\allowbreak{}record\_\allowbreak{}id & bigint u & N & PK & auto-increment \\
vehicle\_\allowbreak{}id & bigint u & N & idx & FK to vehicles.vehicle\_\allowbreak{}id, ON DELETE CASCADE \\
assessment\_\allowbreak{}year & year & N &  &  \\
amount\_\allowbreak{}due & decimal(10,2) & N &  &  \\
due\_\allowbreak{}date & date & N &  &  \\
status & enum('DUE',\allowbreak{}'PAID',\allowbreak{}'OVERDUE') & N &  & default DUE \\
\end{longtable}
\vspace{-0.4em}
\subsection*{Permits (11\_permits)}
\begin{longtable}{@{}p{3.4cm}p{2.9cm}c c p{5.6cm}@{}}
\caption*{\textbf{routes}}\\
\toprule \textbf{Column} & \textbf{Type} & \textbf{Null} & \textbf{Key} & \textbf{Notes} \\ \midrule \endfirsthead
\toprule \textbf{Column} & \textbf{Type} & \textbf{Null} & \textbf{Key} & \textbf{Notes} \\ \midrule \endhead
\bottomrule \endlastfoot
route\_\allowbreak{}id & bigint u & N & PK & auto-increment \\
route\_\allowbreak{}name & varchar(120) & N &  &  \\
origin & varchar(80) & N &  &  \\
destination & varchar(80) & N &  &  \\
\end{longtable}
\vspace{-0.4em}
\begin{longtable}{@{}p{3.4cm}p{2.9cm}c c p{5.6cm}@{}}
\caption*{\textbf{route\_\allowbreak{}segments}}\\
\toprule \textbf{Column} & \textbf{Type} & \textbf{Null} & \textbf{Key} & \textbf{Notes} \\ \midrule \endfirsthead
\toprule \textbf{Column} & \textbf{Type} & \textbf{Null} & \textbf{Key} & \textbf{Notes} \\ \midrule \endhead
\bottomrule \endlastfoot
segment\_\allowbreak{}id & bigint u & N & PK & auto-increment \\
route\_\allowbreak{}id & bigint u & N & idx & FK to routes.route\_\allowbreak{}id, ON DELETE CASCADE \\
sequence\_\allowbreak{}no & smallint u & N &  &  \\
segment\_\allowbreak{}name & varchar(120) & N &  &  \\
\end{longtable}
\vspace{-0.4em}
\begin{longtable}{@{}p{3.4cm}p{2.9cm}c c p{5.6cm}@{}}
\caption*{\textbf{permits}}\\
\toprule \textbf{Column} & \textbf{Type} & \textbf{Null} & \textbf{Key} & \textbf{Notes} \\ \midrule \endfirsthead
\toprule \textbf{Column} & \textbf{Type} & \textbf{Null} & \textbf{Key} & \textbf{Notes} \\ \midrule \endhead
\bottomrule \endlastfoot
permit\_\allowbreak{}id & bigint u & N & PK & auto-increment \\
permit\_\allowbreak{}number & varchar(30) & N & UK &  \\
vehicle\_\allowbreak{}id & bigint u & N & idx & FK to vehicles.vehicle\_\allowbreak{}id, ON DELETE RESTRICT \\
citizen\_\allowbreak{}id & bigint u & N & idx & FK to citizens.citizen\_\allowbreak{}id, ON DELETE RESTRICT \\
permit\_\allowbreak{}type\_\allowbreak{}id & bigint u & N & idx & FK to permit\_\allowbreak{}types.permit\_\allowbreak{}type\_\allowbreak{}id, ON DELETE RESTRICT \\
route\_\allowbreak{}id & bigint u & Y & idx & FK to routes.route\_\allowbreak{}id, ON DELETE SET NULL \\
application\_\allowbreak{}id & bigint u & N & idx & FK to applications.application\_\allowbreak{}id, ON DELETE RESTRICT \\
issue\_\allowbreak{}date & date & N &  &  \\
expiry\_\allowbreak{}date & date & N &  &  \\
status & enum('ACTIVE',\allowbreak{}'EXPIRED',\allowbreak{}'SUSPENDED',\allowbreak{}'CANCELLED') & N &  & default ACTIVE \\
\end{longtable}
\vspace{-0.4em}
\begin{longtable}{@{}p{3.4cm}p{2.9cm}c c p{5.6cm}@{}}
\caption*{\textbf{permit\_\allowbreak{}status\_\allowbreak{}history}}\\
\toprule \textbf{Column} & \textbf{Type} & \textbf{Null} & \textbf{Key} & \textbf{Notes} \\ \midrule \endfirsthead
\toprule \textbf{Column} & \textbf{Type} & \textbf{Null} & \textbf{Key} & \textbf{Notes} \\ \midrule \endhead
\bottomrule \endlastfoot
history\_\allowbreak{}id & bigint u & N & PK & auto-increment \\
permit\_\allowbreak{}id & bigint u & N & idx & FK to permits.permit\_\allowbreak{}id, ON DELETE CASCADE \\
previous\_\allowbreak{}status & varchar(20) & Y &  &  \\
new\_\allowbreak{}status & varchar(20) & N &  &  \\
changed\_\allowbreak{}at & timestamp & N &  & default CURRENT\_\allowbreak{}TIMESTAMP \\
reason & varchar(255) & Y &  &  \\
\end{longtable}
\vspace{-0.4em}
\subsection*{Violations (12\_violations)}
\begin{longtable}{@{}p{3.4cm}p{2.9cm}c c p{5.6cm}@{}}
\caption*{\textbf{violations}}\\
\toprule \textbf{Column} & \textbf{Type} & \textbf{Null} & \textbf{Key} & \textbf{Notes} \\ \midrule \endfirsthead
\toprule \textbf{Column} & \textbf{Type} & \textbf{Null} & \textbf{Key} & \textbf{Notes} \\ \midrule \endhead
\bottomrule \endlastfoot
violation\_\allowbreak{}id & bigint u & N & PK & auto-increment \\
vehicle\_\allowbreak{}id & bigint u & N & idx & FK to vehicles.vehicle\_\allowbreak{}id, ON DELETE RESTRICT \\
driver\_\allowbreak{}citizen\_\allowbreak{}id & bigint u & Y & idx & FK to citizens.citizen\_\allowbreak{}id, ON DELETE SET NULL \\
violation\_\allowbreak{}type\_\allowbreak{}id & bigint u & N & idx & FK to violation\_\allowbreak{}types.violation\_\allowbreak{}type\_\allowbreak{}id, ON DELETE RESTRICT \\
officer\_\allowbreak{}employee\_\allowbreak{}id & bigint u & N & idx & FK to employees.employee\_\allowbreak{}id, ON DELETE RESTRICT \\
location & varchar(200) & N &  &  \\
occurred\_\allowbreak{}at & datetime & N &  &  \\
\end{longtable}
\vspace{-0.4em}
\begin{longtable}{@{}p{3.4cm}p{2.9cm}c c p{5.6cm}@{}}
\caption*{\textbf{challans}}\\
\toprule \textbf{Column} & \textbf{Type} & \textbf{Null} & \textbf{Key} & \textbf{Notes} \\ \midrule \endfirsthead
\toprule \textbf{Column} & \textbf{Type} & \textbf{Null} & \textbf{Key} & \textbf{Notes} \\ \midrule \endhead
\bottomrule \endlastfoot
challan\_\allowbreak{}id & bigint u & N & PK & auto-increment \\
challan\_\allowbreak{}number & varchar(30) & N & UK &  \\
total\_\allowbreak{}amount & decimal(10,2) & N &  &  \\
status & enum('ISSUED',\allowbreak{}'DISPUTED',\allowbreak{}'PAID',\allowbreak{}'CANCELLED') & N & idx & default ISSUED \\
issued\_\allowbreak{}at & datetime & N &  & default CURRENT\_\allowbreak{}TIMESTAMP \\
\end{longtable}
\vspace{-0.4em}
\begin{longtable}{@{}p{3.4cm}p{2.9cm}c c p{5.6cm}@{}}
\caption*{\textbf{challan\_\allowbreak{}violations}}\\
\toprule \textbf{Column} & \textbf{Type} & \textbf{Null} & \textbf{Key} & \textbf{Notes} \\ \midrule \endfirsthead
\toprule \textbf{Column} & \textbf{Type} & \textbf{Null} & \textbf{Key} & \textbf{Notes} \\ \midrule \endhead
\bottomrule \endlastfoot
challan\_\allowbreak{}id & bigint u & N & PK & FK to challans.challan\_\allowbreak{}id, ON DELETE CASCADE \\
violation\_\allowbreak{}id & bigint u & N & PK & FK to violations.violation\_\allowbreak{}id, ON DELETE RESTRICT \\
\end{longtable}
\vspace{-0.4em}
\begin{longtable}{@{}p{3.4cm}p{2.9cm}c c p{5.6cm}@{}}
\caption*{\textbf{challan\_\allowbreak{}status\_\allowbreak{}history}}\\
\toprule \textbf{Column} & \textbf{Type} & \textbf{Null} & \textbf{Key} & \textbf{Notes} \\ \midrule \endfirsthead
\toprule \textbf{Column} & \textbf{Type} & \textbf{Null} & \textbf{Key} & \textbf{Notes} \\ \midrule \endhead
\bottomrule \endlastfoot
history\_\allowbreak{}id & bigint u & N & PK & auto-increment \\
challan\_\allowbreak{}id & bigint u & N & idx & FK to challans.challan\_\allowbreak{}id, ON DELETE CASCADE \\
previous\_\allowbreak{}status & varchar(20) & Y &  &  \\
new\_\allowbreak{}status & varchar(20) & N &  &  \\
changed\_\allowbreak{}at & timestamp & N &  & default CURRENT\_\allowbreak{}TIMESTAMP \\
\end{longtable}
\vspace{-0.4em}
\subsection*{Payments (13\_payments)}
\begin{longtable}{@{}p{3.4cm}p{2.9cm}c c p{5.6cm}@{}}
\caption*{\textbf{fee\_\allowbreak{}structures}}\\
\toprule \textbf{Column} & \textbf{Type} & \textbf{Null} & \textbf{Key} & \textbf{Notes} \\ \midrule \endfirsthead
\toprule \textbf{Column} & \textbf{Type} & \textbf{Null} & \textbf{Key} & \textbf{Notes} \\ \midrule \endhead
\bottomrule \endlastfoot
fee\_\allowbreak{}id & bigint u & N & PK & auto-increment \\
service\_\allowbreak{}type\_\allowbreak{}id & bigint u & N & idx & FK to service\_\allowbreak{}types.service\_\allowbreak{}type\_\allowbreak{}id, ON DELETE CASCADE \\
component\_\allowbreak{}name & varchar(80) & N &  &  \\
amount & decimal(10,2) & N &  &  \\
effective\_\allowbreak{}from & date & N &  &  \\
effective\_\allowbreak{}to & date & Y &  &  \\
\end{longtable}
\vspace{-0.4em}
\begin{longtable}{@{}p{3.4cm}p{2.9cm}c c p{5.6cm}@{}}
\caption*{\textbf{payments}}\\
\toprule \textbf{Column} & \textbf{Type} & \textbf{Null} & \textbf{Key} & \textbf{Notes} \\ \midrule \endfirsthead
\toprule \textbf{Column} & \textbf{Type} & \textbf{Null} & \textbf{Key} & \textbf{Notes} \\ \midrule \endhead
\bottomrule \endlastfoot
payment\_\allowbreak{}id & bigint u & N & PK & auto-increment \\
receipt\_\allowbreak{}number & varchar(30) & N & UK &  \\
payable\_\allowbreak{}type\_\allowbreak{}id & bigint u & N & idx & FK to payable\_\allowbreak{}types.payable\_\allowbreak{}type\_\allowbreak{}id, ON DELETE RESTRICT \\
payable\_\allowbreak{}id & bigint u & N &  &  \\
amount & decimal(10,2) & N &  &  \\
status & enum('PENDING',\allowbreak{}'SUCCESS',\allowbreak{}'FAILED',\allowbreak{}'REFUNDED') & N & idx & default PENDING \\
paid\_\allowbreak{}at & datetime & Y &  &  \\
created\_\allowbreak{}at & datetime & N &  & default CURRENT\_\allowbreak{}TIMESTAMP \\
\end{longtable}
\vspace{-0.4em}
\begin{longtable}{@{}p{3.4cm}p{2.9cm}c c p{5.6cm}@{}}
\caption*{\textbf{payment\_\allowbreak{}attempts}}\\
\toprule \textbf{Column} & \textbf{Type} & \textbf{Null} & \textbf{Key} & \textbf{Notes} \\ \midrule \endfirsthead
\toprule \textbf{Column} & \textbf{Type} & \textbf{Null} & \textbf{Key} & \textbf{Notes} \\ \midrule \endhead
\bottomrule \endlastfoot
attempt\_\allowbreak{}id & bigint u & N & PK & auto-increment \\
payment\_\allowbreak{}id & bigint u & N & idx & FK to payments.payment\_\allowbreak{}id, ON DELETE CASCADE \\
attempt\_\allowbreak{}number & smallint u & N &  &  \\
gateway\_\allowbreak{}reference & varchar(60) & Y &  &  \\
outcome & enum('SUCCESS',\allowbreak{}'FAILED',\allowbreak{}'TIMEOUT') & N &  &  \\
attempted\_\allowbreak{}at & datetime & N &  & default CURRENT\_\allowbreak{}TIMESTAMP \\
\end{longtable}
\vspace{-0.4em}
\begin{longtable}{@{}p{3.4cm}p{2.9cm}c c p{5.6cm}@{}}
\caption*{\textbf{refunds}}\\
\toprule \textbf{Column} & \textbf{Type} & \textbf{Null} & \textbf{Key} & \textbf{Notes} \\ \midrule \endfirsthead
\toprule \textbf{Column} & \textbf{Type} & \textbf{Null} & \textbf{Key} & \textbf{Notes} \\ \midrule \endhead
\bottomrule \endlastfoot
refund\_\allowbreak{}id & bigint u & N & PK & auto-increment \\
payment\_\allowbreak{}id & bigint u & N & idx & FK to payments.payment\_\allowbreak{}id, ON DELETE RESTRICT \\
amount & decimal(10,2) & N &  &  \\
reason & varchar(255) & N &  &  \\
status & enum('INITIATED',\allowbreak{}'COMPLETED',\allowbreak{}'FAILED') & N &  & default INITIATED \\
processed\_\allowbreak{}at & datetime & Y &  &  \\
\end{longtable}
\vspace{-0.4em}
\subsection*{Support (14\_support)}
\begin{longtable}{@{}p{3.4cm}p{2.9cm}c c p{5.6cm}@{}}
\caption*{\textbf{notifications}}\\
\toprule \textbf{Column} & \textbf{Type} & \textbf{Null} & \textbf{Key} & \textbf{Notes} \\ \midrule \endfirsthead
\toprule \textbf{Column} & \textbf{Type} & \textbf{Null} & \textbf{Key} & \textbf{Notes} \\ \midrule \endhead
\bottomrule \endlastfoot
notification\_\allowbreak{}id & bigint u & N & PK & auto-increment \\
person\_\allowbreak{}id & bigint u & N & idx & FK to persons.person\_\allowbreak{}id, ON DELETE CASCADE \\
channel & enum('SMS',\allowbreak{}'EMAIL') & N &  &  \\
subject & varchar(150) & Y &  &  \\
message & varchar(500) & N &  &  \\
status & enum('QUEUED',\allowbreak{}'SENT',\allowbreak{}'FAILED') & N &  & default QUEUED \\
sent\_\allowbreak{}at & datetime & Y &  &  \\
\end{longtable}
\vspace{-0.4em}
\begin{longtable}{@{}p{3.4cm}p{2.9cm}c c p{5.6cm}@{}}
\caption*{\textbf{complaints}}\\
\toprule \textbf{Column} & \textbf{Type} & \textbf{Null} & \textbf{Key} & \textbf{Notes} \\ \midrule \endfirsthead
\toprule \textbf{Column} & \textbf{Type} & \textbf{Null} & \textbf{Key} & \textbf{Notes} \\ \midrule \endhead
\bottomrule \endlastfoot
complaint\_\allowbreak{}id & bigint u & N & PK & auto-increment \\
citizen\_\allowbreak{}id & bigint u & N & idx & FK to citizens.citizen\_\allowbreak{}id, ON DELETE RESTRICT \\
office\_\allowbreak{}id & bigint u & N & idx & FK to rto\_\allowbreak{}offices.office\_\allowbreak{}id, ON DELETE RESTRICT \\
subject & varchar(150) & N &  &  \\
description & varchar(1000) & N &  &  \\
status & enum('OPEN',\allowbreak{}'IN\_\allowbreak{}PROGRESS',\allowbreak{}'RESOLVED',\allowbreak{}'CLOSED') & N &  & default OPEN \\
filed\_\allowbreak{}at & datetime & N &  & default CURRENT\_\allowbreak{}TIMESTAMP \\
\end{longtable}
\vspace{-0.4em}
\begin{longtable}{@{}p{3.4cm}p{2.9cm}c c p{5.6cm}@{}}
\caption*{\textbf{appeals}}\\
\toprule \textbf{Column} & \textbf{Type} & \textbf{Null} & \textbf{Key} & \textbf{Notes} \\ \midrule \endfirsthead
\toprule \textbf{Column} & \textbf{Type} & \textbf{Null} & \textbf{Key} & \textbf{Notes} \\ \midrule \endhead
\bottomrule \endlastfoot
appeal\_\allowbreak{}id & bigint u & N & PK & auto-increment \\
citizen\_\allowbreak{}id & bigint u & N & idx & FK to citizens.citizen\_\allowbreak{}id, ON DELETE RESTRICT \\
against\_\allowbreak{}type & enum('CHALLAN',\allowbreak{}'APPLICATION\_\allowbreak{}REJECTION',\allowbreak{}'LICENCE\_\allowbreak{}SUSPENSION') & N &  &  \\
against\_\allowbreak{}id & bigint u & N &  &  \\
grounds & varchar(1000) & N &  &  \\
status & enum('FILED',\allowbreak{}'UNDER\_\allowbreak{}REVIEW',\allowbreak{}'UPHELD',\allowbreak{}'DISMISSED') & N &  & default FILED \\
filed\_\allowbreak{}at & datetime & N &  & default CURRENT\_\allowbreak{}TIMESTAMP \\
resolved\_\allowbreak{}at & datetime & Y &  &  \\
\end{longtable}
\vspace{-0.4em}
\begin{longtable}{@{}p{3.4cm}p{2.9cm}c c p{5.6cm}@{}}
\caption*{\textbf{audit\_\allowbreak{}logs}}\\
\toprule \textbf{Column} & \textbf{Type} & \textbf{Null} & \textbf{Key} & \textbf{Notes} \\ \midrule \endfirsthead
\toprule \textbf{Column} & \textbf{Type} & \textbf{Null} & \textbf{Key} & \textbf{Notes} \\ \midrule \endhead
\bottomrule \endlastfoot
audit\_\allowbreak{}id & bigint u & N & PK & auto-increment \\
table\_\allowbreak{}name & varchar(64) & N & idx &  \\
record\_\allowbreak{}id & bigint u & N &  &  \\
action & enum('INSERT',\allowbreak{}'UPDATE',\allowbreak{}'DELETE') & N &  &  \\
old\_\allowbreak{}values & json & Y &  &  \\
new\_\allowbreak{}values & json & Y &  &  \\
changed\_\allowbreak{}by\_\allowbreak{}user\_\allowbreak{}id & bigint u & Y &  &  \\
changed\_\allowbreak{}at & timestamp & N & idx & default CURRENT\_\allowbreak{}TIMESTAMP \\
\end{longtable}
\vspace{-0.4em}
\endgroup

\chapter{REST API Operations}\label{app:api}
The backend exposes 225 operations on 151 paths. They were read from the running
application's own API description, so the list matches the code. ``Signed in (service
rules)'' means that any signed-in user may call the operation, and the service itself
decides what that user may see.

\begingroup\footnotesize
\begin{longtable}{@{}p{2.1cm}l p{6.6cm} p{3.2cm}@{}}
\caption{All REST operations exposed by the backend (paths are relative to 	exttt{/\allowbreak{}api/\allowbreak{}v1}), with the permission each requires.}\label{tab:apilist}\\
\toprule \textbf{Group} & \textbf{Verb} & \textbf{Path} & \textbf{Permission} \\ \midrule \endfirsthead
\toprule \textbf{Group} & \textbf{Verb} & \textbf{Path} & \textbf{Permission} \\ \midrule \endhead
\bottomrule \endlastfoot
Application & POST & \texttt{/\allowbreak{}applications} & application.create or application.create\_\allowbreak{}own \\
 & GET & \texttt{/\allowbreak{}applications} & signed in (service rules) \\
 & GET & \texttt{/\allowbreak{}applications/\allowbreak{}\{applicationId\}} & signed in (service rules) \\
 & PATCH & \texttt{/\allowbreak{}applications/\allowbreak{}\{applicationId\}} & application.update \\
 & POST & \texttt{/\allowbreak{}applications/\allowbreak{}\{applicationId\}/\allowbreak{}assign} & application.assign \\
 & POST & \texttt{/\allowbreak{}applications/\allowbreak{}\{applicationId\}/\allowbreak{}status} & signed in (service rules) \\
 & GET & \texttt{/\allowbreak{}applications/\allowbreak{}\{applicationId\}/\allowbreak{}history} & signed in (service rules) \\
 & POST & \texttt{/\allowbreak{}applications/\allowbreak{}\{applicationId\}/\allowbreak{}documents} & document.upload \\
 & GET & \texttt{/\allowbreak{}applications/\allowbreak{}\{applicationId\}/\allowbreak{}documents} & signed in (service rules) \\
 & GET & \texttt{/\allowbreak{}documents/\allowbreak{}\{documentId\}/\allowbreak{}file} & signed in (service rules) \\
 & PATCH & \texttt{/\allowbreak{}documents/\allowbreak{}\{documentId\}/\allowbreak{}verify} & document.verify \\
 & PATCH & \texttt{/\allowbreak{}documents/\allowbreak{}\{documentId\}/\allowbreak{}reject} & document.verify \\
 & PATCH & \texttt{/\allowbreak{}documents/\allowbreak{}\{documentId\}/\allowbreak{}expire} & document.verify \\
Appointment & GET & \texttt{/\allowbreak{}appointment-slots} & signed in (service rules) \\
 & POST & \texttt{/\allowbreak{}appointment-slots} & appointment.manage \\
 & PATCH & \texttt{/\allowbreak{}appointment-slots/\allowbreak{}\{slotId\}} & appointment.manage \\
 & POST & \texttt{/\allowbreak{}applications/\allowbreak{}\{applicationId\}/\allowbreak{}appointments} & appointment.book \\
 & GET & \texttt{/\allowbreak{}appointments} & signed in (service rules) \\
 & PATCH & \texttt{/\allowbreak{}appointments/\allowbreak{}\{appointmentId\}/\allowbreak{}cancel} & appointment.book \\
 & PATCH & \texttt{/\allowbreak{}appointments/\allowbreak{}\{appointmentId\}/\allowbreak{}reschedule} & appointment.book \\
 & PATCH & \texttt{/\allowbreak{}appointments/\allowbreak{}\{appointmentId\}/\allowbreak{}outcome} & appointment.manage \\
Auth & POST & \texttt{/\allowbreak{}auth/\allowbreak{}login} & public \\
 & POST & \texttt{/\allowbreak{}auth/\allowbreak{}refresh} & public \\
 & POST & \texttt{/\allowbreak{}auth/\allowbreak{}logout} & signed in (service rules) \\
 & GET & \texttt{/\allowbreak{}auth/\allowbreak{}me} & signed in (service rules) \\
Citizen & POST & \texttt{/\allowbreak{}citizens} & citizen.create \\
 & GET & \texttt{/\allowbreak{}citizens} & citizen.view \\
 & GET & \texttt{/\allowbreak{}citizens/\allowbreak{}me} & signed in (service rules) \\
 & GET & \texttt{/\allowbreak{}citizens/\allowbreak{}\{citizenId\}} & signed in (service rules) \\
 & PATCH & \texttt{/\allowbreak{}citizens/\allowbreak{}\{citizenId\}} & citizen.update \\
 & GET & \texttt{/\allowbreak{}citizens/\allowbreak{}\{citizenId\}/\allowbreak{}addresses} & signed in (service rules) \\
 & POST & \texttt{/\allowbreak{}citizens/\allowbreak{}\{citizenId\}/\allowbreak{}addresses} & citizen.update \\
Compliance & GET & \texttt{/\allowbreak{}vehicles/\allowbreak{}\{vehicleId\}/\allowbreak{}compliance} & compliance.view \\
 & GET & \texttt{/\allowbreak{}compliance/\allowbreak{}expiring} & compliance.view \\
 & POST & \texttt{/\allowbreak{}compliance/\allowbreak{}refresh-statuses} & compliance.manage \\
 & POST & \texttt{/\allowbreak{}inspections} & compliance.manage \\
 & GET & \texttt{/\allowbreak{}inspections} & compliance.view \\
 & GET & \texttt{/\allowbreak{}inspections/\allowbreak{}\{inspectionId\}} & compliance.view \\
 & POST & \texttt{/\allowbreak{}fitness-certificates} & compliance.manage \\
 & GET & \texttt{/\allowbreak{}fitness-certificates} & compliance.view \\
 & GET & \texttt{/\allowbreak{}fitness-certificates/\allowbreak{}\{certificateId\}} & compliance.view \\
 & POST & \texttt{/\allowbreak{}fitness-certificates/\allowbreak{}\{certificateId\}/\allowbreak{}revoke} & compliance.manage \\
 & POST & \texttt{/\allowbreak{}pollution-certificates} & compliance.manage \\
 & GET & \texttt{/\allowbreak{}pollution-certificates} & compliance.view \\
 & GET & \texttt{/\allowbreak{}pollution-certificates/\allowbreak{}\{pucId\}} & compliance.view \\
 & POST & \texttt{/\allowbreak{}insurance-policies} & compliance.manage \\
 & GET & \texttt{/\allowbreak{}insurance-policies} & compliance.view \\
 & GET & \texttt{/\allowbreak{}insurance-policies/\allowbreak{}\{policyId\}} & compliance.view \\
 & POST & \texttt{/\allowbreak{}road-tax} & compliance.manage \\
 & GET & \texttt{/\allowbreak{}road-tax} & compliance.view \\
 & GET & \texttt{/\allowbreak{}road-tax/\allowbreak{}\{taxRecordId\}} & compliance.view \\
Employee & POST & \texttt{/\allowbreak{}employees} & employee.manage \\
 & GET & \texttt{/\allowbreak{}employees} & employee.view \\
 & GET & \texttt{/\allowbreak{}employees/\allowbreak{}\{employeeId\}} & employee.view \\
 & PATCH & \texttt{/\allowbreak{}employees/\allowbreak{}\{employeeId\}} & employee.manage \\
 & GET & \texttt{/\allowbreak{}employees/\allowbreak{}\{employeeId\}/\allowbreak{}postings} & employee.view \\
 & POST & \texttt{/\allowbreak{}employees/\allowbreak{}\{employeeId\}/\allowbreak{}postings} & employee.manage \\
Licence & POST & \texttt{/\allowbreak{}driving-schools} & licence.manage\_\allowbreak{}schools \\
 & GET & \texttt{/\allowbreak{}driving-schools} & licence.view \\
 & GET & \texttt{/\allowbreak{}driving-schools/\allowbreak{}\{schoolId\}} & licence.view \\
 & PATCH & \texttt{/\allowbreak{}driving-schools/\allowbreak{}\{schoolId\}} & licence.manage\_\allowbreak{}schools \\
 & POST & \texttt{/\allowbreak{}driving-schools/\allowbreak{}\{schoolId\}/\allowbreak{}instructors} & licence.manage\_\allowbreak{}schools \\
 & GET & \texttt{/\allowbreak{}driving-schools/\allowbreak{}\{schoolId\}/\allowbreak{}instructors} & licence.view \\
 & POST & \texttt{/\allowbreak{}test-centres} & licence.manage\_\allowbreak{}schools \\
 & GET & \texttt{/\allowbreak{}test-centres} & licence.view \\
 & POST & \texttt{/\allowbreak{}learner-licences} & licence.issue \\
 & GET & \texttt{/\allowbreak{}learner-licences} & licence.view \\
 & GET & \texttt{/\allowbreak{}learner-licences/\allowbreak{}\{learnerId\}} & signed in (service rules) \\
 & POST & \texttt{/\allowbreak{}learner-licences/\allowbreak{}\{learnerId\}/\allowbreak{}cancel} & licence.revoke \\
 & POST & \texttt{/\allowbreak{}driving-licences} & licence.issue \\
 & GET & \texttt{/\allowbreak{}driving-licences} & licence.view \\
 & POST & \texttt{/\allowbreak{}driving-licences/\allowbreak{}expire-overdue} & licence.revoke \\
 & GET & \texttt{/\allowbreak{}driving-licences/\allowbreak{}\{licenceId\}} & signed in (service rules) \\
 & POST & \texttt{/\allowbreak{}driving-licences/\allowbreak{}\{licenceId\}/\allowbreak{}classes} & licence.issue \\
 & POST & \texttt{/\allowbreak{}driving-licences/\allowbreak{}\{licenceId\}/\allowbreak{}suspend} & licence.suspend \\
 & POST & \texttt{/\allowbreak{}driving-licences/\allowbreak{}\{licenceId\}/\allowbreak{}revoke} & licence.revoke \\
 & POST & \texttt{/\allowbreak{}driving-licences/\allowbreak{}\{licenceId\}/\allowbreak{}reinstate} & licence.issue \\
 & POST & \texttt{/\allowbreak{}driving-licences/\allowbreak{}\{licenceId\}/\allowbreak{}renew} & licence.issue \\
 & GET & \texttt{/\allowbreak{}driving-licences/\allowbreak{}\{licenceId\}/\allowbreak{}history} & signed in (service rules) \\
 & POST & \texttt{/\allowbreak{}driving-tests} & licence.test \\
 & GET & \texttt{/\allowbreak{}driving-tests} & licence.view \\
 & GET & \texttt{/\allowbreak{}driving-tests/\allowbreak{}\{testId\}} & licence.view \\
 & PATCH & \texttt{/\allowbreak{}driving-tests/\allowbreak{}\{testId\}/\allowbreak{}result} & licence.test \\
Office & POST & \texttt{/\allowbreak{}offices} & office.manage \\
 & GET & \texttt{/\allowbreak{}offices} & signed in (service rules) \\
 & GET & \texttt{/\allowbreak{}offices/\allowbreak{}\{officeId\}} & signed in (service rules) \\
 & PATCH & \texttt{/\allowbreak{}offices/\allowbreak{}\{officeId\}} & office.manage \\
 & GET & \texttt{/\allowbreak{}offices/\allowbreak{}\{officeId\}/\allowbreak{}counters} & signed in (service rules) \\
 & POST & \texttt{/\allowbreak{}offices/\allowbreak{}\{officeId\}/\allowbreak{}counters} & office.manage \\
 & PATCH & \texttt{/\allowbreak{}counters/\allowbreak{}\{counterId\}} & office.manage \\
 & DELETE & \texttt{/\allowbreak{}counters/\allowbreak{}\{counterId\}} & office.manage \\
 & GET & \texttt{/\allowbreak{}departments} & signed in (service rules) \\
 & POST & \texttt{/\allowbreak{}departments} & office.manage \\
 & GET & \texttt{/\allowbreak{}designations} & signed in (service rules) \\
 & POST & \texttt{/\allowbreak{}designations} & employee.manage \\
 & PATCH & \texttt{/\allowbreak{}designations/\allowbreak{}\{designationId\}} & employee.manage \\
Payment & POST & \texttt{/\allowbreak{}payments} & payment.create \\
 & GET & \texttt{/\allowbreak{}payments} & payment.view \\
 & GET & \texttt{/\allowbreak{}payments/\allowbreak{}\{paymentId\}} & payment.view \\
 & GET & \texttt{/\allowbreak{}payments/\allowbreak{}\{paymentId\}/\allowbreak{}attempts} & payment.view \\
 & POST & \texttt{/\allowbreak{}payments/\allowbreak{}\{paymentId\}/\allowbreak{}retry} & payment.create \\
 & POST & \texttt{/\allowbreak{}payments/\allowbreak{}\{paymentId\}/\allowbreak{}refund} & payment.refund \\
 & GET & \texttt{/\allowbreak{}payments/\allowbreak{}\{paymentId\}/\allowbreak{}refunds} & payment.view \\
 & GET & \texttt{/\allowbreak{}refunds} & payment.view \\
 & GET & \texttt{/\allowbreak{}refunds/\allowbreak{}\{refundId\}} & payment.view \\
 & POST & \texttt{/\allowbreak{}refunds/\allowbreak{}\{refundId\}/\allowbreak{}complete} & payment.refund \\
 & POST & \texttt{/\allowbreak{}refunds/\allowbreak{}\{refundId\}/\allowbreak{}fail} & payment.refund \\
 & POST & \texttt{/\allowbreak{}challans/\allowbreak{}\{challanId\}/\allowbreak{}pay} & payment.create \\
 & GET & \texttt{/\allowbreak{}fee-structures} & signed in (service rules) \\
 & POST & \texttt{/\allowbreak{}fee-structures} & reference.manage \\
 & PATCH & \texttt{/\allowbreak{}fee-structures/\allowbreak{}\{feeId\}} & reference.manage \\
Permit & POST & \texttt{/\allowbreak{}routes} & route.manage \\
 & GET & \texttt{/\allowbreak{}routes} & permit.view \\
 & GET & \texttt{/\allowbreak{}routes/\allowbreak{}\{routeId\}} & permit.view \\
 & POST & \texttt{/\allowbreak{}routes/\allowbreak{}\{routeId\}/\allowbreak{}segments} & route.manage \\
 & GET & \texttt{/\allowbreak{}routes/\allowbreak{}\{routeId\}/\allowbreak{}segments} & permit.view \\
 & POST & \texttt{/\allowbreak{}permits} & permit.create \\
 & GET & \texttt{/\allowbreak{}permits} & permit.view \\
 & POST & \texttt{/\allowbreak{}permits/\allowbreak{}expire-overdue} & permit.manage \\
 & GET & \texttt{/\allowbreak{}permits/\allowbreak{}\{permitId\}} & permit.view \\
 & GET & \texttt{/\allowbreak{}permits/\allowbreak{}\{permitId\}/\allowbreak{}history} & permit.view \\
 & POST & \texttt{/\allowbreak{}permits/\allowbreak{}\{permitId\}/\allowbreak{}suspend} & permit.manage \\
 & POST & \texttt{/\allowbreak{}permits/\allowbreak{}\{permitId\}/\allowbreak{}reinstate} & permit.manage \\
 & POST & \texttt{/\allowbreak{}permits/\allowbreak{}\{permitId\}/\allowbreak{}cancel} & permit.manage \\
Rbac & GET & \texttt{/\allowbreak{}permissions} & rbac.manage \\
 & GET & \texttt{/\allowbreak{}roles} & rbac.manage \\
 & POST & \texttt{/\allowbreak{}roles} & rbac.manage \\
 & GET & \texttt{/\allowbreak{}roles/\allowbreak{}\{roleId\}} & rbac.manage \\
 & PUT & \texttt{/\allowbreak{}roles/\allowbreak{}\{roleId\}/\allowbreak{}permissions} & rbac.manage \\
 & GET & \texttt{/\allowbreak{}users} & user.view \\
 & POST & \texttt{/\allowbreak{}users} & user.manage \\
 & GET & \texttt{/\allowbreak{}users/\allowbreak{}\{userId\}} & user.view \\
 & PATCH & \texttt{/\allowbreak{}users/\allowbreak{}\{userId\}} & user.manage \\
 & PUT & \texttt{/\allowbreak{}users/\allowbreak{}\{userId\}/\allowbreak{}roles} & rbac.manage \\
Reference & GET & \texttt{/\allowbreak{}regions} & signed in (service rules) \\
 & GET & \texttt{/\allowbreak{}regions/\allowbreak{}\{itemId\}} & signed in (service rules) \\
 & POST & \texttt{/\allowbreak{}regions} & reference.manage \\
 & PATCH & \texttt{/\allowbreak{}regions/\allowbreak{}\{itemId\}} & reference.manage \\
 & DELETE & \texttt{/\allowbreak{}regions/\allowbreak{}\{itemId\}} & reference.manage \\
 & GET & \texttt{/\allowbreak{}vehicle-manufacturers} & signed in (service rules) \\
 & GET & \texttt{/\allowbreak{}vehicle-manufacturers/\allowbreak{}\{itemId\}} & signed in (service rules) \\
 & POST & \texttt{/\allowbreak{}vehicle-manufacturers} & reference.manage \\
 & PATCH & \texttt{/\allowbreak{}vehicle-manufacturers/\allowbreak{}\{itemId\}} & reference.manage \\
 & DELETE & \texttt{/\allowbreak{}vehicle-manufacturers/\allowbreak{}\{itemId\}} & reference.manage \\
 & GET & \texttt{/\allowbreak{}vehicle-models} & signed in (service rules) \\
 & GET & \texttt{/\allowbreak{}vehicle-models/\allowbreak{}\{itemId\}} & signed in (service rules) \\
 & POST & \texttt{/\allowbreak{}vehicle-models} & reference.manage \\
 & PATCH & \texttt{/\allowbreak{}vehicle-models/\allowbreak{}\{itemId\}} & reference.manage \\
 & DELETE & \texttt{/\allowbreak{}vehicle-models/\allowbreak{}\{itemId\}} & reference.manage \\
 & GET & \texttt{/\allowbreak{}vehicle-types} & signed in (service rules) \\
 & GET & \texttt{/\allowbreak{}vehicle-types/\allowbreak{}\{itemId\}} & signed in (service rules) \\
 & POST & \texttt{/\allowbreak{}vehicle-types} & reference.manage \\
 & PATCH & \texttt{/\allowbreak{}vehicle-types/\allowbreak{}\{itemId\}} & reference.manage \\
 & DELETE & \texttt{/\allowbreak{}vehicle-types/\allowbreak{}\{itemId\}} & reference.manage \\
 & GET & \texttt{/\allowbreak{}fuel-types} & signed in (service rules) \\
 & GET & \texttt{/\allowbreak{}fuel-types/\allowbreak{}\{itemId\}} & signed in (service rules) \\
 & POST & \texttt{/\allowbreak{}fuel-types} & reference.manage \\
 & PATCH & \texttt{/\allowbreak{}fuel-types/\allowbreak{}\{itemId\}} & reference.manage \\
 & DELETE & \texttt{/\allowbreak{}fuel-types/\allowbreak{}\{itemId\}} & reference.manage \\
 & GET & \texttt{/\allowbreak{}licence-classes} & signed in (service rules) \\
 & GET & \texttt{/\allowbreak{}licence-classes/\allowbreak{}\{itemId\}} & signed in (service rules) \\
 & POST & \texttt{/\allowbreak{}licence-classes} & reference.manage \\
 & PATCH & \texttt{/\allowbreak{}licence-classes/\allowbreak{}\{itemId\}} & reference.manage \\
 & DELETE & \texttt{/\allowbreak{}licence-classes/\allowbreak{}\{itemId\}} & reference.manage \\
 & GET & \texttt{/\allowbreak{}document-types} & signed in (service rules) \\
 & GET & \texttt{/\allowbreak{}document-types/\allowbreak{}\{itemId\}} & signed in (service rules) \\
 & POST & \texttt{/\allowbreak{}document-types} & reference.manage \\
 & PATCH & \texttt{/\allowbreak{}document-types/\allowbreak{}\{itemId\}} & reference.manage \\
 & DELETE & \texttt{/\allowbreak{}document-types/\allowbreak{}\{itemId\}} & reference.manage \\
 & GET & \texttt{/\allowbreak{}service-types} & signed in (service rules) \\
 & GET & \texttt{/\allowbreak{}service-types/\allowbreak{}\{itemId\}} & signed in (service rules) \\
 & POST & \texttt{/\allowbreak{}service-types} & reference.manage \\
 & PATCH & \texttt{/\allowbreak{}service-types/\allowbreak{}\{itemId\}} & reference.manage \\
 & DELETE & \texttt{/\allowbreak{}service-types/\allowbreak{}\{itemId\}} & reference.manage \\
 & GET & \texttt{/\allowbreak{}violation-types} & signed in (service rules) \\
 & GET & \texttt{/\allowbreak{}violation-types/\allowbreak{}\{itemId\}} & signed in (service rules) \\
 & POST & \texttt{/\allowbreak{}violation-types} & reference.manage \\
 & PATCH & \texttt{/\allowbreak{}violation-types/\allowbreak{}\{itemId\}} & reference.manage \\
 & DELETE & \texttt{/\allowbreak{}violation-types/\allowbreak{}\{itemId\}} & reference.manage \\
 & GET & \texttt{/\allowbreak{}permit-types} & signed in (service rules) \\
 & GET & \texttt{/\allowbreak{}permit-types/\allowbreak{}\{itemId\}} & signed in (service rules) \\
 & POST & \texttt{/\allowbreak{}permit-types} & reference.manage \\
 & PATCH & \texttt{/\allowbreak{}permit-types/\allowbreak{}\{itemId\}} & reference.manage \\
 & DELETE & \texttt{/\allowbreak{}permit-types/\allowbreak{}\{itemId\}} & reference.manage \\
 & GET & \texttt{/\allowbreak{}regions/\allowbreak{}tree} & signed in (service rules) \\
 & GET & \texttt{/\allowbreak{}payable-types} & signed in (service rules) \\
Support & POST & \texttt{/\allowbreak{}complaints} & complaint.create \\
 & GET & \texttt{/\allowbreak{}complaints} & signed in (service rules) \\
 & GET & \texttt{/\allowbreak{}complaints/\allowbreak{}\{complaintId\}} & signed in (service rules) \\
 & PATCH & \texttt{/\allowbreak{}complaints/\allowbreak{}\{complaintId\}/\allowbreak{}status} & complaint.manage \\
 & POST & \texttt{/\allowbreak{}appeals} & appeal.create \\
 & GET & \texttt{/\allowbreak{}appeals} & signed in (service rules) \\
 & GET & \texttt{/\allowbreak{}appeals/\allowbreak{}\{appealId\}} & signed in (service rules) \\
 & POST & \texttt{/\allowbreak{}appeals/\allowbreak{}\{appealId\}/\allowbreak{}review} & appeal.review \\
 & POST & \texttt{/\allowbreak{}appeals/\allowbreak{}\{appealId\}/\allowbreak{}uphold} & appeal.review \\
 & POST & \texttt{/\allowbreak{}appeals/\allowbreak{}\{appealId\}/\allowbreak{}dismiss} & appeal.review \\
 & GET & \texttt{/\allowbreak{}notifications} & signed in (service rules) \\
 & POST & \texttt{/\allowbreak{}notifications/\allowbreak{}expiry-reminders} & licence.issue or notification.view \\
 & GET & \texttt{/\allowbreak{}audit-logs} & audit.view \\
 & GET & \texttt{/\allowbreak{}dashboard/\allowbreak{}summary} & dashboard.view \\
System & GET & \texttt{/\allowbreak{}health} & public \\
Vehicle & POST & \texttt{/\allowbreak{}vehicles} & vehicle.create \\
 & GET & \texttt{/\allowbreak{}vehicles} & vehicle.view \\
 & GET & \texttt{/\allowbreak{}vehicles/\allowbreak{}\{vehicleId\}} & vehicle.view \\
 & PATCH & \texttt{/\allowbreak{}vehicles/\allowbreak{}\{vehicleId\}} & vehicle.update \\
 & GET & \texttt{/\allowbreak{}vehicles/\allowbreak{}\{vehicleId\}/\allowbreak{}owners} & vehicle.view \\
 & GET & \texttt{/\allowbreak{}vehicles/\allowbreak{}\{vehicleId\}/\allowbreak{}current-owner} & vehicle.view \\
 & POST & \texttt{/\allowbreak{}vehicles/\allowbreak{}\{vehicleId\}/\allowbreak{}owners} & vehicle.update \\
 & GET & \texttt{/\allowbreak{}citizens/\allowbreak{}\{citizenId\}/\allowbreak{}vehicles} & signed in (service rules) \\
 & POST & \texttt{/\allowbreak{}ownership-transfers} & vehicle.transfer \\
 & GET & \texttt{/\allowbreak{}ownership-transfers} & vehicle.view \\
 & GET & \texttt{/\allowbreak{}ownership-transfers/\allowbreak{}\{transferId\}} & vehicle.view \\
 & POST & \texttt{/\allowbreak{}ownership-transfers/\allowbreak{}\{transferId\}/\allowbreak{}approve} & vehicle.transfer\_\allowbreak{}approve \\
 & POST & \texttt{/\allowbreak{}ownership-transfers/\allowbreak{}\{transferId\}/\allowbreak{}reject} & vehicle.transfer\_\allowbreak{}approve \\
Violation & POST & \texttt{/\allowbreak{}violations} & violation.create \\
 & GET & \texttt{/\allowbreak{}violations} & violation.view \\
 & GET & \texttt{/\allowbreak{}violations/\allowbreak{}\{violationId\}} & violation.view \\
 & POST & \texttt{/\allowbreak{}challans} & challan.create \\
 & GET & \texttt{/\allowbreak{}challans} & challan.view \\
 & GET & \texttt{/\allowbreak{}challans/\allowbreak{}\{challanId\}} & challan.view \\
 & GET & \texttt{/\allowbreak{}challans/\allowbreak{}\{challanId\}/\allowbreak{}history} & challan.view \\
 & POST & \texttt{/\allowbreak{}challans/\allowbreak{}\{challanId\}/\allowbreak{}violations} & challan.manage \\
 & POST & \texttt{/\allowbreak{}challans/\allowbreak{}\{challanId\}/\allowbreak{}dispute} & challan.manage \\
 & POST & \texttt{/\allowbreak{}challans/\allowbreak{}\{challanId\}/\allowbreak{}reissue} & challan.manage \\
 & POST & \texttt{/\allowbreak{}challans/\allowbreak{}\{challanId\}/\allowbreak{}cancel} & challan.manage \\
\end{longtable}
\endgroup

\chapter{Permission Matrix}\label{app:perms}
The default grants of the seven roles. Role names are abbreviated as ADM (administrator),
OFF (RTO officer), CLK (clerk), INS (inspector), ENF (enforcement officer), ACC (accounts
officer) and CIT (citizen). The grants are data in the database and can be changed from the
role screen of the portal (Figure~\ref{fig:scrrole}).

\begingroup\footnotesize
\begin{longtable}{@{}p{4.6cm}ccccccc@{}}
\caption{Default permission grants of the seven seeded roles (a dot means the role holds the permission).}\label{tab:permmatrix}\\
\toprule \textbf{Permission} & \textbf{ADM} & \textbf{OFF} & \textbf{CLK} & \textbf{INS} & \textbf{ENF} & \textbf{ACC} & \textbf{CIT} \\ \midrule \endfirsthead
\toprule \textbf{Permission} & \textbf{ADM} & \textbf{OFF} & \textbf{CLK} & \textbf{INS} & \textbf{ENF} & \textbf{ACC} & \textbf{CIT} \\ \midrule \endhead
\bottomrule \endlastfoot
rbac.manage & $\bullet$ &  &  &  &  &  &  \\
user.view & $\bullet$ &  &  &  &  &  &  \\
user.manage & $\bullet$ &  &  &  &  &  &  \\
audit.view & $\bullet$ &  &  &  &  & $\bullet$ &  \\
dashboard.view & $\bullet$ & $\bullet$ & $\bullet$ &  &  & $\bullet$ &  \\
reference.manage & $\bullet$ &  &  &  &  &  &  \\
citizen.view & $\bullet$ & $\bullet$ & $\bullet$ & $\bullet$ & $\bullet$ &  &  \\
citizen.create & $\bullet$ & $\bullet$ & $\bullet$ &  &  &  &  \\
citizen.update & $\bullet$ & $\bullet$ & $\bullet$ &  &  &  &  \\
citizen.view\_\allowbreak{}sensitive & $\bullet$ &  &  &  &  &  &  \\
office.manage & $\bullet$ &  &  &  &  &  &  \\
employee.view & $\bullet$ & $\bullet$ &  &  &  &  &  \\
employee.manage & $\bullet$ &  &  &  &  &  &  \\
application.view & $\bullet$ & $\bullet$ & $\bullet$ &  &  & $\bullet$ &  \\
application.view\_\allowbreak{}own & $\bullet$ &  &  &  &  &  & $\bullet$ \\
application.create & $\bullet$ & $\bullet$ & $\bullet$ &  &  &  &  \\
application.create\_\allowbreak{}own & $\bullet$ &  &  &  &  &  & $\bullet$ \\
application.update & $\bullet$ & $\bullet$ & $\bullet$ &  &  &  &  \\
application.assign & $\bullet$ & $\bullet$ &  &  &  &  &  \\
application.verify & $\bullet$ & $\bullet$ &  &  &  &  &  \\
application.approve & $\bullet$ & $\bullet$ &  &  &  &  &  \\
application.reject & $\bullet$ & $\bullet$ &  &  &  &  &  \\
application.cancel & $\bullet$ & $\bullet$ &  &  &  &  &  \\
document.upload & $\bullet$ & $\bullet$ & $\bullet$ &  &  &  & $\bullet$ \\
document.view & $\bullet$ & $\bullet$ & $\bullet$ &  &  &  &  \\
document.verify & $\bullet$ & $\bullet$ &  &  &  &  &  \\
appointment.view & $\bullet$ & $\bullet$ & $\bullet$ &  &  &  &  \\
appointment.manage & $\bullet$ & $\bullet$ & $\bullet$ &  &  &  &  \\
appointment.book & $\bullet$ & $\bullet$ & $\bullet$ &  &  &  & $\bullet$ \\
licence.view & $\bullet$ & $\bullet$ & $\bullet$ & $\bullet$ & $\bullet$ &  &  \\
licence.issue & $\bullet$ & $\bullet$ &  &  &  &  &  \\
licence.suspend & $\bullet$ & $\bullet$ &  &  &  &  &  \\
licence.revoke & $\bullet$ & $\bullet$ &  &  &  &  &  \\
licence.test & $\bullet$ & $\bullet$ &  & $\bullet$ &  &  &  \\
licence.manage\_\allowbreak{}schools & $\bullet$ &  &  &  &  &  &  \\
vehicle.view & $\bullet$ & $\bullet$ & $\bullet$ & $\bullet$ & $\bullet$ &  &  \\
vehicle.create & $\bullet$ & $\bullet$ &  &  &  &  &  \\
vehicle.update & $\bullet$ & $\bullet$ &  &  &  &  &  \\
vehicle.transfer & $\bullet$ & $\bullet$ & $\bullet$ &  &  &  &  \\
vehicle.transfer\_\allowbreak{}approve & $\bullet$ & $\bullet$ &  &  &  &  &  \\
compliance.view & $\bullet$ & $\bullet$ & $\bullet$ & $\bullet$ &  &  &  \\
compliance.manage & $\bullet$ & $\bullet$ &  & $\bullet$ &  &  &  \\
permit.view & $\bullet$ & $\bullet$ & $\bullet$ &  &  &  &  \\
permit.create & $\bullet$ & $\bullet$ &  &  &  &  &  \\
permit.manage & $\bullet$ & $\bullet$ &  &  &  &  &  \\
route.manage & $\bullet$ & $\bullet$ &  &  &  &  &  \\
violation.view & $\bullet$ & $\bullet$ &  & $\bullet$ & $\bullet$ &  &  \\
violation.create & $\bullet$ &  &  & $\bullet$ & $\bullet$ &  &  \\
challan.view & $\bullet$ & $\bullet$ & $\bullet$ &  & $\bullet$ & $\bullet$ &  \\
challan.create & $\bullet$ &  &  &  & $\bullet$ &  &  \\
challan.manage & $\bullet$ &  &  &  & $\bullet$ &  &  \\
payment.view & $\bullet$ & $\bullet$ & $\bullet$ &  &  & $\bullet$ &  \\
payment.create & $\bullet$ &  & $\bullet$ &  &  & $\bullet$ &  \\
payment.refund & $\bullet$ &  &  &  &  & $\bullet$ &  \\
complaint.create & $\bullet$ &  & $\bullet$ &  &  &  & $\bullet$ \\
complaint.view\_\allowbreak{}own & $\bullet$ &  &  &  &  &  & $\bullet$ \\
complaint.manage & $\bullet$ & $\bullet$ &  &  &  &  &  \\
appeal.create & $\bullet$ &  &  &  &  &  & $\bullet$ \\
appeal.view\_\allowbreak{}own & $\bullet$ &  &  &  &  &  & $\bullet$ \\
appeal.review & $\bullet$ & $\bullet$ &  &  &  &  &  \\
notification.view & $\bullet$ & $\bullet$ &  &  &  &  &  \\
notification.view\_\allowbreak{}own & $\bullet$ &  &  &  &  &  & $\bullet$ \\
\midrule \textbf{Total} & 62 & 41 & 21 & 8 & 8 & 7 & 9 \\
\end{longtable}
\endgroup

\chapter{Build and Run Guide}\label{app:run}
\section{Prerequisites}
JDK 21, Maven 3.9, MySQL 8.0 and a modern browser. The portal loads its style library from
a content-delivery network, so an internet connection is needed the first time a page is
opened.

\section{Steps}
\begin{enumerate}
\item \textbf{Create the database.} In the \cd{database} folder, run the script \cd{00_run_all.sql} with the MySQL client as a user who may create databases. It creates \cd{rto_management} and its 63 tables.
\item \textbf{Create an application user} with all privileges on \cd{rto_management} (and on \cd{rto_management_test} if the tests will be run).
\item \textbf{Configure the backend.} Copy \cd{backend-java/.env.example} to \cd{backend-java/.env} and fill in the database host, port, user and password, a signing secret of at least 32 characters, and the name and password of the first administrator. This file must never be committed.
\item \textbf{Build and prepare.} In \cd{backend-java}, run \cmd{mvn -DskipTests package}, then run the jar once with \cmd{--rto.command=bootstrap} to load the roles, permissions and the first administrator.
\item \textbf{Optional demonstration data.} Load the comprehensive seed from \cd{database/seed} so that every screen has records to show.
\item \textbf{Start the backend} (port 8000) with \texttt{java -jar target/\allowbreak{}rto-management-\allowbreak{}backend-1.0.0.jar}. The interactive API documentation is served at \cd{/docs}.
\item \textbf{Build and start the portal} (port 8080) from \cd{frontend-java} in the same way. It reaches the backend through the setting \cd{RTO_API_BASE_URL}, which defaults to the local backend.
\item \textbf{Sign in} at \cd{http://localhost:8080} with the administrator account.
\item \textbf{Run the tests} with \cmd{mvn test} in \cd{backend-java}. They drop and rebuild the test database.
\end{enumerate}

\section{Repeating the Measurements}
The load test, the index experiment and the crawl used short scripts that are not part of
the application. Their numerical results are tabulated in Chapters~\ref{ch:perf}
and~\ref{ch:test}, and the method of each is described there so that it can be repeated.
